%%%%%%%%%%%%%%%%%%%%%%%%%%%%%%%%%%%%%%%%%%%%%%%%%%%%%%%%%%%%%%%%%%%%%%%%%%%%%%%%
%%                           ONDER DE LOEP ARTIKEL                            %%
%%%%%%%%%%%%%%%%%%%%%%%%%%%%%%%%%%%%%%%%%%%%%%%%%%%%%%%%%%%%%%%%%%%%%%%%%%%%%%%%

\artikelfoto{Negen fragmenten uit de geschiedenis van de wiskunde}{}{Michel Roelens \\ Luc Van den Broeck \\ Els Vanlommel}{img/LoepGeschiedenis_coverfoto 3903 kleiner.jpg}

\startcontents[inhoudloep]	% om inhoud voor het loep artikel te beginnen

\begin{multicols}{2}

    \inhoudstafelloep
	
	\section{Inleiding}
	 Deze loep bestaat uit negen historische fragmenten uit de geschiedenis van de wiskunde die passen bij bepaalde leerstofonderdelen. Sommige fragmenten kun je de avond voor de les nog lezen. Andere geschiedkundige tussendoortjes vragen wat meer voorbereidingstijd. Je kunt korte, geselecteerde stukjes van deze loep gebruiken naargelang de gelegenheid en de interesse van de leerlingen zich voordoet. Niets van deze loep hoort bij de verplichte leerstof. Maar zoals je al eerder in de klas zal ervaren hebben: wat niet verplicht is, is dubbel zo interessant.

    Een idee-fixe bij vele mensen is dat de wiskunde altijd al geweest is wat ze nu is, dat ze al eeuwenlang quasi onveranderlijk werd doorgegeven, net zoals de bijbel. Het beschrijven van de soms moeizame evolutie van de wiskunde doorheen verschillende tijden en culturen is belangrijk om te kunnen begrijpen waar we nu mee bezig zijn. 
    
    Uiteraard zullen we het in deze loep hebben over de periode van de klassieke oudheid toen de getallen nog niet waren wat ze nu zijn en de irrationale getallen bijvoorbeeld niet meer waren dan 'onderlinge onmeetbaarheid van lijnstukken'. Maar we zullen het ook hebben over de wiskunde van de periode na 1500, toen de negatieve getallen en de decimale kommagetallen pas in gebruik begonnen te raken. Dat het zoveel eeuwen duurde vooraleer de westerse wereld het getalinzicht verwierf dat wij nu als elementair bestempelen, noopt ons tot nederigheid en misschien ook tot bewondering voor al wat nadien nog ontdekt is. 

    Omdat de meeste wiskunde uit het secundair onderwijs al eeuwen bekend is, hebben we het hier niet over de hedendaagse wiskunde. Aan kopstukken zoals Euler en Gauss kunnen we uiteraard niet voorbij gaan. Recentere wiskunde zoals de 23 problemen voor de twintigste eeuw van Hilbert, de Fieldsmedaille en de Belgische medaillewinnaars, de Bourbakisten en de moderne wiskunde, de doorbraak in het bewijs van het vermoeden van Fermat, nemen we niet op. Hiervoor hebben we geen marge. 

    In elk van de negen fragmenten vermelden we duidelijk bij welke leerstof en in welke graad de uitweiding past. Soms is er keuze tussen verschillende graden en tussen verschillende leerstofonderdelen. 
    
    De uitwerking van de verschillende fragmenten kan variëren. Sommige stukjes focussen op het \emph{redeneren}. We proberen hierin klassieke bewijzen vanuit de vertaalde teksten te begrijpen en te reconstrueren. In deze passages vind je ingekaderde lesactiviteiten, die individueel of in groepjes of met de klas kunnen uitgevoerd worden. Om de werkteksten wat toegankelijker te maken, hebben we oude teksten voorzien van moderne notaties: functionele anachronismen ten voordele van de begrijpbaarheid.
    
    Andere stukjes focussen op de \emph{historische context} van de wiskunde. Een bijna teloor gegane activiteit in de wiskundeles is het vertellen. In onze jonge tijd was het niet ongewoon dat een wiskundeleerkracht op regelmatige tijdstippen achter zijn lessenaar gingen zitten om schijnbaar onvoorbereid bepaalde topics uit de geschiedenis van de wiskunde uit zijn mouw te schudden. We noteerden niets maar hingen aan zijn lippen. Ook dit lesverloop willen we voorzichtig aanmoedigen. Lees op avond voor de les een brokje geschiedenis uit deze loep en onthou er enkele fragmenten van. Wedden dat de leerlingen van je verhaal genieten?

    Naast het \emph{redeneren} en het gebruiken van de \emph{historische context} kunnen we de geschiedenis van de wiskunde in de klas ook gebruiken door middel van \emph{anekdotiek}. Ongeacht het geromantiseerde karakter van deze verhalen, zijn het net de anekdoten die bij elke leerling blijven hangen. Het zijn ook de kapstokken waaraan andere wiskundige informatie wordt opgehangen. Welke leerling smult er nu niet van het eureka-verhaal van Archimedes of van zijn uitspraak ``verstoor mijn cirkels niet'' net voor hij bij de inname van Syracuse naast zijn zandbak met wiskundige schetsen gedood werd door een Romeinse soldaat? Wie bejegent Gauss nu niet met gepaste egards wanneer hij het verhaal voorgeschoteld krijgt van de jonge Carl Friedrich die zijn meester in de lagere school te slim af was toen hij de getallen van 1 tot 100 moest optellen?

    We zijn er ons bewust van dat deze loep slechts een kleine selectie is uit het aanbod aan geschiedkundige passages die hun diensten kunnen bewijzen in onze lessen. Voor andere en dieper uitgewerkte fragmenten, verwijzen we naar het boek \emph{Wortels van de Wiskunde. Een historisch overzicht voor leraren en anderen}. Hierin vind je de geschiedenis van de wiskunde in een grote notendop samen met 25 uitgelezen hoofdstukken over specifieke onderwerpen.

    Tot slot van deze inleiding, als je geïnteresseerd bent in een grondigere uitwerking van een bepaald thema uit de geschiedenis van de wiskunde, grasduin dan eerst even met de zoekrobot door onze website. We werkten vroeger al een loep uit over \emph{Christiaan Huyghens als gastdocent} (UW 10/1), over \emph{Het astrolabium} (29/4), over \emph{Veeltermvergelijkingen van vroeger tot nu} (UW 31/4) enz. 
	
	
	\section{Simon Stevin en de decimale getallen}
	\subsection*{Waarom deze keuze?}
    
    Doorheen heel het secundair onderwijs worden verschillende getalverzamelingen bestudeerd: natuurlijke getallen, negatieve getallen, breuken, decimale getallen en complexe getallen. De meeste getalverzamelingen zijn zo vanzelfspekend dat het er op lijkt dat wetenschappers ze altijd al hebben gekend en gebruikt. 
    
    De leerlingen zijn verbaasd als ze vernemen dat het getal 0, oorspronkelijk een Indisch concept, pas rond 1100 na Christus in het Westen is doorgedrongen door Leonardo van Pisa (alias Fibonacci), die het tijdens zijn veelvuldige reizen door Noord-Afrika overnam van de Arabieren. Nog meer verbazing is er over de negatieve getallen. Toen Scipio del Ferro zich rond 1515 verdiepte in de oplossing van derdegraadsvergelijkingen, kon hij enkel steunen op positieve getallen. De negatieve waren toen nog niet in gebruik. 
    
    In het begin van de zestiende eeuw was er ook nog geen notatie voor decimale kommagetallen. Een eerste aanzet kwam er  in 1585. De naar Leiden uitgeweken Bruggeling Simon Stevin publiceerde toen een van zijn belangrijkste werken, \emph{De Thiende}, waarin hij het tiendelig positiestelsel aanbeval. Er zou nog heel wat tijd overgaan eer de decimale getallen eruit zagen zoals wij ze gewoon zijn. Simon Stevin schreef de `decimale kommagetallen' nog zonder komma. Nadien was het nog wachten op Napoleon vooraleer lengtematen, inhoudsmaten en gewichtsmaten rond 1800 in brede kringen gestandaardiseerd werden in het tiendelige stelsel.

    Simon Stevin kan terecht de Archimedes van de Lage Landen genoemd worden. Hij is een wetenschapper, een uitvinder, een taalkundige, een waterbouwkundige, een stad- en vestingbouwer, ... Samen met Mercator en Vesalius belichaamde hij de modernisering van de wetenschap in de zestiende eeuw. Daarom werd hij onlangs opgenomen in de Canon van Vlaanderen. Voor de West-Vlamingen kan hij beschouwd worden als een `streekgenoot'. Ook daarom is het zinvol een beetje tijd voor hem uit te trekken in onze lessen. 

    \afbeelding[6cm]{img/LoepGeschiedenis_stevin.jpg}{Simon Stevin (1548-1620)}{fig:stevin}

    \subsection*{Bij welke gelegenheid?}
    
    Thematisch past dit praatje over Stevin in de eerste en de tweede graad wanneer we het hebben over de reële getallen. We kunnen hem ook vermelden wanneer we het hebben over de wiskundige terminologie die we vandaag gebruiken. Verschillende namen van begrippen stammen nog uit de tijd waarin het Latijn de voertaal voor wetenschappers was. Maar bijna alle wiskundige begripsnamen met een Nederlandse herkomst zijn door Simon Stevin bedacht.
    

    \subsection*{Het leven van Simon Stevin}

    Simon Stevin werd in 1548 in Brugge geboren als onwettig kind van Antheunis Stevin en Cathelyne van der Poort. Van Stevins jeugd is weinig geweten. Zijn openbare leven begon pas bij zijn uitwijking naar Leiden in 1571. Rond die tijd was Brugge een broeihaard van conflicten tussen de calvinisten en de katholieken. Het was ook de tijd van de beeldenstormen, waarin calvinistische iconoclasten zich afzetten tegen de Spaanse bezetter door de kunstschatten van katholieke kerken te vernielen. Simon Stevin was vermoedelijk protestant. Na zijn studie aan de vrijzinnige universiteit van Leiden kon hij zich opwerken tot de vertrouwenspersoon van Prins Maurits, stadhouder van Holland en Zeeland.

    In 1593 stelde Maurits van Oranje, zoon van Willem van Oranje, Stevin aan als privéleraar. Mauritz was immers erg geïnteresseerd in wetenschappen en in de militaire en waterbouwkundige toepassingen ervan. De inhoud van deze lessen werd gebundeld in de \emph{Wisconstighe Ghedachtenissen}. Stevin werd levenslang onderhouden door zijn mecenas, Prins Mauritz, die hij in de \emph{Wisconstighe Ghedachtenissen} met alle eer overlaadde. Stevin kreeg van hem een functie als kwartiermeester in het Staatse Leger en hij mocht naast de universiteit van Leiden een Nederlandstalige ingenieursopleiding uitwerken met klemtonen op wiskunde, landmeetkunde, waterbeheersing en vestingbouw.
      
    In 1600, midden in de tachtigjarige oorlog, probeerde prins Mauritz de troepen van de Spaanse overheerser te verdrijven uit Duinkerkerke, een bolwerk van Spaanse kapers. Dit leidde tot de beruchte slag bij Nieuwpoort, die nipt door de Staatse Troepen van Mauritz gewonnen werd. Mauritz kon nadien wel een oppeppertje gebruiken. Hiertoe, en ook ten voordele van zijn eigen populariteit, ontwierp Stevin een zeilwagen waarmee hij in het voorjaar van 1602 samen met een aantal notabelen een rit op het strand maakte van Scheveningen naar Petten. Op deze 90 km werd een gemiddelde snelheid van 40 km/h gehaald. Ter vergelijking, momenteel ligt het snelheidsrecord van een strandzeilwagen boven de 200 km/h.

    \afbeelding[8cm]{img/LoepGeschiedenis_zeilwagen.jpg}{Zeilwagen op het strand van Scheveningen, mei 1602}{fig:zeilwagen}
 
    In 1586 publiceerde Stevin drie wetenschappelijke werken: \emph{De beghinselen der weeghconst}, \emph{De weeghdaet} en \emph{De beghinselen des waterwichts}. In het eerste werk beschrijft hij hoe hij samen met een bevriende geleerde, Jan Cornets, een valproef doet vanop de meer dan 100 meter hoge toren van de Nieuwe Kerk in Delft. Vandaar liet hij twee loden ballen (in het Nederlands van toen: loyen cloten, cloot is verwant met kluit) met een verschillende massa op een houten vloer naast de toren vallen. De ballen vertrokken samen en kwamen samen aan. Hierdoor toonde Stevin nog voor Galileo Galilei aan dat de valsnelheid niet afhangt van de massa van het vallende lichaam. Galileo ging echter onterecht aan de haal met de eer van de eerste valproef.
    \subsection*{De Thiende}

    In de tijd van Simon Stevin was er een wildgroei aan maataanduidingen die niet tiendelig waren. Er waren zestigdelige, zoals ons systeem van minuten en seconden, maar men rekende ook twaalfdelig met dozijnen en grossen. In Groot-Brittanië rekende men tot 1971 nog met ponden die onderverdeeld waren in 20 shillings en schillings die verdeeld waren in 12 pence. Andere niet-decimale maten in Groot-Brittanië zijn de yard die bestaat uit 3 feet en de foot die bestaat uit 12 inches. In Amerika zijn het de gallon, de quart, de pint, ... 
    
    In 1585 schreef Simon Stevin een klein maar belangrijk werkje, \emph{De Thiende}, waarin hij de overstap naar een tiendelige maatstelsel aanraadde. Stevin draagt dit werk op aan de gewone burgers: landmeters, tapijtmeters, muntmeesters en kooplieden. Hijzelf zorgde voor een Franse vertaling \emph{La disme}. Snel volgde er een Engelse en een Deense vertaling.

    \renewcommand*\circled[1]{\tikz[baseline=(char.base)]{
            \node[shape=circle,draw,inner sep=2pt] (char) {#1};}}
            
    In dit werkje stelde hij ook een nieuwe notatie voor decimale breuken voor, die handiger was dan de notatie met breukstrepen met noemers 10, 100, 1000 enz. Hij plaatste \circled{0} achter het cijfer van de eenheden, \circled{1} achter het cijfer van de tienden, \circled{2} achter de honderdtallen ... Het getal 7,564 noteerde Stevin dus als 
    \[7 \; \circled{0} \; 5 \; \circled{1} \; 6 \; \circled{2} \; 4 \; \circled{3}\;.\]
    Hij introduceerde met andere woorden een tiendelig positiestelsel voor breuken waarbij de positie van elk cijfer nog eens extra benadrukt wordt met een omcirkeld getal.

    Het grote voordeel van deze notatie is dat er makkelijk kan gecijferd worden. Stevin werkt enkele voorbeelden uit in \emph{De Thiende} en verklaart ze ook. De omcirkelde cijfers staan in deze schema's niet \emph{in} de getallen, maar \emph{erboven} of \emph{eronder}. Wij herkennen in deze voorbeelden makkelijk de cijferoefeningen die we in de lagere school maakten.

    \afbeelding[4cm]{img/LoepGeschiedenis_som.jpg}{De som van 27,847 en 37,675 en 875,782}{fig:som}

    \afbeelding[5.5cm]{img/LoepGeschiedenis_product.jpg}{Het product van 32,57 en 89,46}{fig:product}

    \end{multicols}
    \clearpage

    \afbeelding[12cm]{img/stevin.jpg}{}{}{}

    \begin{multicols}{2}

    Stevin was terecht heel fier op deze notatie. In zijn werk verwoordt hij dit als volgt:
    
    \begin{citaat}
    Als iemand mij derhalve voor een opschepper zou houden enkel en alleen omdat ik het nut van deze methode duidelijk maak, dan toont hij daarmee aan ofwel dat hij geen benul heeft van het onderscheid tussen het gewone en het bijzondere, dan wel dat hij het algemeen belang een slecht hart toedraagt.
    \end{citaat}

 
    
    
    \subsection*{Het Nederlands als wetenschapstaal}

    Het Latijn was de wetenschapstaal bij uitstek in de zestiende eeuw. Nu is het Engels de belangrijkste internationale wetenschapstaal. Stevin zette in twee verhandelingen argumenten uit om aan te tonen dat het Nederlands een nuttigere voertaal zou zijn dan het Latijn. Hij merkte vooreerst op dat het Nederlands wendbaarder is door de overvloed aan eenlettergrepige woorden. Om dit aan te tonen, stelde Stevin een lange lijst samen van eenlettergrepige woorden en hun equivalenten in het Frans en het Latijn.
    \begin{citaat}
        Caert. Ie ioue aux cartes. Ludo chartis.\\
        Stroom. Cours de l'eau. Fluxus aquæ \\
        Schaeck. Ie prens vne fille par force. Rapio virginem.
    \end{citaat}

    Een tweede argument voor de wendbaarheid van het Nederlands was het kunnen samenvoegen van twee woorden tot één waarbij de volgorde van de woorden van belang is voor de betekenis. Zo is putwater niet hetzelfde als waterput en verschilt de jachthond van de hondenjacht en vensterglas van glasvenster.

    Om de wetenschapstaal bij uitstek te worden, stelde Stevin een aantal wetenschapstermen voor die de Latijnse moesten vervangen. Sommige waren al beperkt in gebruik zoals wortel (voor radix) en drijhoeck (voor triangulus). Andere voorstellen waren nieuw zoals omtreck (voor peripheria), veelhouck (voor polygonum), stelling (voor propositio), raaklijn (voor tangens linea), meetconst (voor geometria) en wisconst (voor mathematica). Ook de andere woorden die op -const eindigen, komen van hem: scheikunde, natuurkunde, aardrijkskunde, stelkunde, \ldots

    Voorstellen van Stevin die het niet gehaald hebben zijn gront (voor basis), scheefrondt (voor ellipsis), brantsne (voor parabola), naelde (voor pyramide), evewijdich vierhouck (voor parallelogramum) en seul (voor cilinder).

    Het zou echter nog enkele eeuwen duren vooraleer het Latijn in onbruik geraakte in wiskundige kringen. De grote wiskundige Carl Friedrich Gauss (1777-1855) publiceerde zijn laatste werk in het Latijn omstreeks 1835. Daarna schakelde hij over op het Duits.

	\section{Euclides en de priemgetallen} \label{euclides}
    \subsection*{Waarom deze keuze?}
 De priemgetallen hebben altijd iets fascinerends en mysterieus. Ze lijken willekeurig verdeeld over de natuurlijke getallen, maar toch heeft de mensheid hierover grote ontdekkingen gedaan en zijn er modellen voor de `asymptotische' verdeling van heel grote priemgetallen (Hadamard en de la Vallée Poussin, rond 1900). Het laatste woord is hierover nog niet gezegd. Men vermoedt dat er oneindig veel priemtweelingen zijn (priemgetallen die opeenvolgende oneven getallen zijn, zoals 5 en 7 of zoals 17 en 19), maar niemand krijgt dit bewezen. Ook de legendarische riemannhypothese, één van de millenniumproblemen die bij oplossing meteen een miljoen dollar opleveren, houdt verband met de verdeling van de priemgetallen.
 
Wat wel al heel lang met zekerheid geweten is, is dat er oneindig veel priemgetallen bestaan. Het bewijs van Euclides (derde eeuw v.C.) is, zeker in een modernere formulering, eigenlijk best toegankelijk en begrijpbaar, ook voor leerlingen van de tweede graad. Een indrukwekkend resultaat met een redelijk eenvoudig bewijs: reden genoeg om dit eeuwenoude fragment in deze loep op te nemen.

 \subsection*{Bij welke gelegenheid?}
In de tweede graad zien de leerlingen het bewijs van de irrationaliteit van $\sqrt{2}$, een bewijs uit het ongerijmde. Misschien is het een goed idee om dit bewijs rond hetzelfde moment te tonen, om het idee van een bewijs uit het ongerijmde vast te zetten. Het vormt ook een toepassing van het hoofdstuk logica in de tweede graad. 

 \subsection*{Wiskundige theorie als een gebouw}
 De Elementen van Euclides (derde eeuw v.C.) zou je de wiskundige encyclopedie van de Oudheid kunnen noemen. In 13 boeken worden de vlakke meetkunde,  ruimtemeetkunde en  getallenleer opgebouwd als één grote constructie, waarbij elke propositie (stelling, eigenschap) bewezen wordt steunend op hetgeen ervoor al bewezen is.
 
 \afbeelding[7cm]{img/LoepGeschiedenis_eerstegedrukteEucl.jpg}{Twee pagina's uit de eerste gedrukte editie van de Elementen}{}
 
Euclides moest hierbij natuurlijk ergens beginnen: een beperkt aantal definities, axioma's en postulaten die als spelregels aangenomen worden en waar al de rest op steunt. Een voorbeeld van een postulaat is: \emph{Door twee punten gaat juist één rechte.}

 \subsection*{Wie was Euclides?}

 Over Euclides van Alexandrië, de auteur van de Elementen, is niet zo veel geweten. Hij was rond 300 v.C. werkzaam in de grote bibliotheek van Alexandrië (nu Egypte), een groot studiecentrum waar wetenschappers uit verschillende landen samenkwamen. De meeste verwijzingen naar Euclides dateren van eeuwen later; niemand heeft hem live geïnterviewd of tijdens zijn leven geportretteerd.

  \subsection*{Het bewijs van Euclides over priemgetallen}

 De Elementen van Euclides is vooral bekend voor de meetkunde, maar bevat ook heel wat getallenleer. In het negende boek vinden we onder meer een bewijs van de stelling dat er oneindig veel priemgetallen zijn.  
 
 In de volgende lesactiviteit maken de leerlingen eerst kennis met een moderne versie van het bewijs van Euclides. Het is dezelfde redenering, maar genoteerd en geformuleerd zoals wij dit nu zouden doen. Op het einde laten we hen ook het eeuwenoude bewijs lezen, in een Nederlandse vertaling. Het enige wat we behalve de taal aangepast hebben, is de notatie voor lijnstukken, om verwarring bij de leerlingen te vermijden. Immers, in de tekst van Euclides worden lijnstukken de ene keer met twee hoofdletters en de andere keer met één hoofdletter aangeduid en we vonden het niet interessant om de leerlingen deze extra hindernis op te leggen.

 Wil je minder tijd steken in deze historische uitweiding, dan kun je je beperken tot de modernere versie van het bewijs.
 Je stopt de lesactiviteit dan na vraag 8.
 
 We veronderstellen dat de leerlingen vooraf even opgefrist hebben wat priemgetallen zijn.

\end{multicols}
\clearpage
 
\begin{lesactiviteit}{Oneindig veel priemgetallen: het bewijs van Euclides}

Met welke redenering kon Euclides bewijzen dat er oneindig veel priemgetallen bestaan? We bekijken zijn bewijs in een moderne versie en in aparte stukjes (kaders). Op het einde laten we de letterlijke tekst van Euclides zien, in vertaling.\\

\begin{kader}{Bewijs}{Er zijn oneindig veel priemgetallen}

Veronderstel dat er maar een eindig aantal priemgetallen $p_1, p_2, ..., p_k$ zijn voor zekere $k\in\mathbb{N}_0$. We gaan bewijzen dat er nog meer priemgetallen zijn dan die.

\end{kader}

\begin{vraagenantwoord}

\vraag{Welk soort bewijs is het?}
\antwoord{Een bewijs uit het ongerijmde.}
\vraag{Wat moet er uiteindelijk uit deze veronderstelling volgen?}
\antwoord{Een tegenspraak. Hier wordt al aangekondigd wat die tegenspraak zal zijn: dat er meer priemgetallen zijn dan die. Dit vormt dan een tegenspraak met de veronderstelling: je veronderstelt dat er geen andere zijn dan die $k$ priemgetallen en er zijn er toch nog meer dan die. De concluisie uit die tegenspraak zal zijn dat de veronderstelling fout is.}

\end{vraagenantwoord}

We gaan verder met het bewijs.\\

\begin{kader}{}{}

Neem het getal $n = p_1\cdot p_2\cdot ...\cdot p_k +1$. Dan is $n$ ofwel priem ofwel niet priem.

\end{kader}

\begin{vraagenantwoord}[resume]

\vraag{Wat is dat getal $n$ als het veronderstelde lijstje van alle priemgetallen $\{2, 3, 5, 7\}$ zou zijn? (We weten natuurlijk dat dit belachelijk kort lijstje niet alle priemgetallen kan bevatten, omdat we nog meer priemgetallen kennen zoals $13$ of $29$, maar we doen dit om de berekening uit het bewijs concreet te maken.)}

\antwoord{Dan is $n$ gelijk aan: $2\cdot3\cdot5\cdot7+1=211$.}

\vraag{Is in dit geval $n$ een priemgetal?}

\antwoord{Ja, $211$ is een priemgetal.}

\vraag{Zelfde vragen voor het lijstje $\{2, 3, 5, 7, 11, 13\}$.}

\antwoord{Dan vinden we $n=2\cdot3\cdot5\cdot7\cdot 11\cdot 13+1=30031$. Dit is geen priemgetal; het is $59\cdot 509$, zoals je online kunt vinden.}
    
\end{vraagenantwoord}

Het bewijs gaat verder\ldots\\

\begin{kader}{}{}

Veronderstel eerst dat $n$ wel priem is. Dan is $n$ een priemgetal dat groter is dan de getallen uit de lijst $p_1, p_2, ..., p_k$. De veronderstelling dat de lijst alle priemgetallen bevat, is dus fout.

Veronderstel nu dat $n$ niet priem is; dan  moet het deelbaar zijn door een zeker priemgetal $p$. Het getal $p$ is niet gelijk aan één van de getallen $p_1, p_2, ..., p_k$. Want, stel dat dit het geval zou zijn, dan zou $p$ een deler zijn zowel van $n=p_1\cdot p_2\cdot ...\cdot p_k+1$ als van $p_1\cdot p_2\cdot ...\cdot p_k$ en dus ook van het verschil $1$, en dit is absurd. Dus is $p$ een priemgetal dat niet in de lijst $p_1, p_2, ..., p_k$ voorkomt. De veronderstelling dat de lijst alle priemgetallen bevat, is dus ook in dit geval fout.

Hiermee is de stelling bewezen.

\end{kader}

\begin{vraagenantwoord}[resume]

\vraag{In beide gevallen wordt een priemgetal gevonden dat nog niet in het lijstje stond. Pas dit toe op de concrete korte lijstjes uit de vragen 3 en 5.}

\antwoord{In het geval van vraag 3 is $n=211$ een priemgetal dat niet in het lijstje stond. In het geval van vraag 5 is $p$ bijvoorbeeld $59$, en dat is ook een priemgetal dat niet in de lijst stond.}

\end{vraagenantwoord}

Binnen dit bewijs uit het ongerijmde staat een mini-bewijs uit het ongerijmde, om uit te leggen dat de priemfactor $p$ niet gelijk kan zijn aan één van de priemgetallen uit de lijst. Hierbij gebruikt men de eigenschap: ``Een deler van twee getallen is ook een deler van het verschil van deze getallen.''

\begin{vraagenantwoord}[resume]
    
\vraag{Bewijs deze eigenschap.}

\antwoord{Veronderstel dat $a$ een deler is van $b$ en van $c$. Dan bestaan er natuurlijke getallen $e$ en $f$ zodat $b=e\cdot a$ en $c=f\cdot a$. Dan is  $b-c=e\cdot a - f\cdot a=(e-f)\cdot a$ en dus is $a$ een deler van $b-c$.}

\vraag{Leg nog eens precies uit hoe deze eigenschap gebruikt wordt om aan te tonen dat $p$ niet gelijk is aan één van de priemgetallen van de lijst.}

\antwoord{Als $p$ gelijk zou zijn aan één van de priemgetallen $p_1, p_2, ..., p_k$, dan zou $p$ een deler zijn van het product $p_1\cdot p_2\cdot ...\cdot p_k$. Maar $p$ is ook een deler van $n=p_1\cdot p_2\cdot ...\cdot p_k+1$. Door de eigenschap zou $p$ dan ook gelijk zijn aan $n$ min dat product. En dat is $1$. En dat kan niet, want geen enkel priemgetal is een deler van 1. Deze tegenspraak leidt tot het verwerpen van de veronderstelling van dit mini-bewijs uit het ongerijmde: $p$ kan niet één van de priemgetallen van de lijst zijn.}

\end{vraagenantwoord}

Op de volgende pagina lees je het bewijs zoals Euclides het geschreven heeft. Het is vrij letterlijk vertaald naar het Nederlands, en de notaties van lijnstukken hebben we ook aangepast aan wat jullie gewoon zijn. Je zult zien dat hij de natuurlijke getallen op een voor ons vreemde manier voorstelt. Ook deelbaarheid bekijkt hij anders dan wij\ldots

Denk bij het bewijs na over de volgende vragen.
\begin{vraagenantwoord}[resume]

\vraag{Hoe stelt hij de getallen voor?}

\antwoord{De natuurlijke getallen worden hier voorgesteld als lijnstukken die veelvouden zijn van `de eenheid'.}

\vraag{Werkt Euclides ook met een willekeurig aantal priemgetallen in de lijst ($k$)?}

\antwoord{Neen, hij bedoelt ook wel een willekeurig aantal maar hij gebruikt in het bewijs een lijstje van slechts drie priemgetallen.}

\vraag{Wat bedoelt Euclides als hij zegt dat getal $[GH]$ `gemeten' kan worden door bv. getal $[AB]$?}

\antwoord{Hij bedoelt dat $[GH]$ een veelvoud is van $[AB]$. Als je het meet met $[AB]$, dit betekent stokjes van lengte $|AB|$ achter elkaar leggen op $[GH]$, komt het juist uit.}

\vraag{Het kleinste getal dat gemeten kan worden door $[AB]$, $[CD]$ en $[EF]$, welk getal is dat?}

\antwoord{Het kleinste gemene veelvoud van die $[AB]$, $[CD]$ en $[EF]$. Maar omdat het priemgetallen zijn, is dat het product van deze drie getallen.}

\end{vraagenantwoord}

\begin{kader}{Boek IX propositie 20 (Euclides, 3de eeuw v.C.)}{Er zijn meer priemgetallen dan elke vooropgestelde hoeveelheid priemgetallen}

\afbeelding[9cm]{img/LoepGeschiedenis_ges0.jpg}{}{ges0}

Gegeven een vooropgestelde hoeveelheid priemgetallen: $[AB]$, $[CD]$ en $[EF]$. Ik beweer dat er nog meer priemgetallen zijn dan $[AB]$, $[CD]$ en $[EF]$. 

Neem het kleinste getal dat gemeten kan worden door $[AB]$, $[CD]$ en $[EF]$; stel dat het $[GH]$ is. Verleng $[GH]$ met de eenheid $[HI]$. Dan is $[GI]$ ofwel priem ofwel niet priem. 

Veronderstel eerst dat $[GI]$ wel priem is; dan hebben we priemgetallen $[AB]$, $[CD]$, $[EF]$ en $[GI]$ gevonden, dus meer dan $[AB]$, $[CD]$ en $[EF]$. 

Veronderstel nu dat $[GI]$ niet priem is; dan moet het meetbaar zijn door een zeker priemgetal (boek VII propositie 31). Neem aan dat het meetbaar is door priemgetal $[JK]$. Ik beweer dat $[JK]$ niet gelijk is aan een van de getallen $[AB]$, $[CD]$ en $[EF]$. Want, stel dat dit het geval zou zijn, dan zou je $[GH]$ kunnen meten met $[JK]$. Maar $[JK]$ meet ook $[GI]$. Dan zou het getal $[JK]$ ook de rest meten, de eenheid $[HI]$, en dit is absurd. Daarom is $[JK]$ niet gelijk aan $[AB]$, $[CD]$ en $[EF]$, maar wel priem. Dus hebben we de priemgetallen $[AB]$, $[CD]$, $[EF]$ en $[JK]$, en dat zijn er meer dan de vooropgestelde $[AB]$, $[CD]$ en $[EF]$ alleen.
\end{kader}

$ $

\end{lesactiviteit}

\begin{multicols}{2}

\subsection*{Bewijs van Saidak}
Ook van wiskundige stellingen die al sinds vele eeuwen bewezen zijn, wordt wel eens een nieuw bewijs gevonden. Misschien is het bewijs van Saidak van 2006 het vermelden waard, om aan leerlingen te laten zien dat wiskunde nog steeds springlevend is.

In tegenstelling tot dat van Euclides, is het geen bewijs uit het ongerijmde. Saidak toont hiernaast rechtstreeks aan dat er voor elke $k\in \mathbb{N}$ minstens $k$ priemgetallen bestaan, zonder evenwel die priemgetallen concreet te construeren.
\vfill\null
\columnbreak

\begin{kader}{Bewijs van Saidak (2006)}{Er zijn oneindig veel priemgetallen}

Stel $n$ is een natuurlijk getal groter dan $1$. De getallen $n$ en $n + 1$ verschillen slechts $1$ en hebben dus geen gemeenschappelijke priemfactoren. Dat betekent dat het getal $n_2 = n(n + 1)$ ten minste twee verschillende priemfactoren heeft. 

Voor de getallen $n_2$ en $n_2 + 1$ geldt hetzelfde: zij verschillen slechts $1$ en hebben dus geen gemeenschappelijke priemfactoren. Het getal $n_3 = n_2(n_2 + 1)$ heeft dus minimaal drie verschillende priemfactoren. 

Als je dit verder zet, zie je dat het getal $n_k$ ten minste $k$ verschillende priemfactoren heeft. Omdat dit voor elk natuurlijk getal $k$ geldt, kan de rij van de priemgetallen nooit ophouden.

\end{kader}
\clearpage

\section{Historische oefeningen over de stelling van Pythagoras}
	 \subsection*{Waarom deze keuze?}

    We hadden een fragment kunnen voorzien over Pythagoras, zijn `school' de pythagoreeërs, hun wereldbeeld en waarom hun `stelling van Pythagoras' dit wereldbeeld aan het wankelen bracht. Bij de aanbreng van de stelling van Pythagoras in de tweede graad is het zeker goed om hier wat meer over te vertellen. Maar daar gaat deze paragraaf niet over. 
    
    Wat we hier wel willen laten zien, is dat er oefeningen zijn uit verschillende periodes en culturen, die met de stelling van Pythagoras kunnen worden opgelost. We denken dat de historische herkomst deze oefeningen extra boeiend kunnen maken. Bovendien geven ze de gelegenheid om iets te vertellen over Chinese en Japanse wiskunde 

    \subsection*{Bij welke gelegenheid?}
  
    In de tweede graad lossen de leerlingen vraagstukken op met de stelling van Pythagoras. Voor mooie oefeningen kunnen we historische inspiratie vinden in het verre Oosten. We bespreken hier een Chinese en een Japanse opgave.

    \subsection*{Het Chinese bamboeprobleem}

    Het wiskundige standaardwerk van het Oude China is de Jiǔzhāng Suànshù (10de tot 2de eeuw v.C.), de Negen Hoofdstukken van de Rekenkunde. Het bevatte algemene oplossingsmethodes voor heel diverse problemen. Het grote verschil met de Griekse traditie is dat er geen bewijzen bij staan. Vele eeuwen later (3de eeuw n.C.) schreef Liú Huī commentaren bij de Negen Hoofdstukken, met verklaringen en bewijzen.

    Het negende van de Negen Hoofdstukken gaat over wat wij de stelling van Pythagoras noemen. De Chinezen noemden die de gōugǔ-stelling, letterlijk: de stelling over de kortere en langere zijden van de rechthoekige driehoek. Eén van de problemen die in dat hoofdstuk opgelost worden, is het bamboeprobleem.

    Je ziet hiernaast een geknakte bamboestengel. In de tekst op de afbeelding wordt het probleem gesteld en opgelost. Voor wie moeite heeft met de Chinese tekst vermelden we de opgave ook in het Nederlands.

    \vfill\null
    \columnbreak
    \afbeelding[5.5cm]{img/LoepGeschiedenis_ges2.JPG}{Het Chinese bamboeprobleem}{ges2}

\begin{kader}{Het Chinese bamboeprobleem}{}

De oorspronkelijke hoogte van de stengel was 10 meter. De top van de stengel is neergekomen op 3 meter afstand van de voet van de stengel. Op welke hoogte is de stengel geknakt? 

\end{kader}

Noem de gevraagde hoogte $x$. De schuine zijde is $10-x$. De gōugǔ-stelling geeft:
\[x^2+3^2=(10-x)^2\]
wat als oplossing \(x=\frac{91}{20}\) geeft.

De Chinezen werkten niet met deze concrete afmetingen (10 meter en 3 meter). Ze gaven een oplossingsrecept in volzinnen om het probleem op te lossen voor willekeurige oorspronkelijke lengte van de stengel en willekeurige horizontale afstand tussen de top en de voet na afknakken: 
\begin{citaat}
    Deel het kwadraat van de horizontale afstand door de lengte, trek dit af van de oorspronkelijk lengte en halveer. 
\end{citaat}
We kunnen hier een formule van maken. Noem $a$ de oorspronkelijke lengte van de stengel en $d$ de horizontale afstand, dan is de hoogte $x$ van wat er van de stengel nog recht staat, volgens dit recept gegeven door:
\[x=\frac{1}{2}\left(a-\frac{d^2}{a}\right)\]

Je kunt controleren dat je hiermee inderdaad $\frac{91}{20}$ vindt als je $a=10$ en $d=3$ invult.

Hoe zijn ze van de stelling van Pythagoras/gōugǔ tot die formule gekomen? Misschien zo (van den Bogaart, 2020):
\[
\begin{array}{lrl}
    &x^2+d^2 &=(a-x)^2 \\
    \Rightarrow &d^2&=(a-x)^2-x^2  \\
    & &=(a-x-x)(a-x+x)\\
    & &=(a-2x)a\\
\end{array}
\]
waar je vervolgens \(x\) uit kunt isoleren:
\[ x =\frac{1}{2}\left(a-\frac{d^2}{a}\right).\]

De Oude Chinezen noteerden dit natuurlijk niet met onze \mbox{algebraïsche} notaties, maar misschien kwam hun redenering op hetzelfde neer.

\subsection*{Een sangaku}

Sangaku's zijn kleurrijke meetkundige puzzels op houten plankjes, die in Japanse tempels hangen. Ze zijn opgebouwd uit cirkels, driehoeken, vierkanten... die mooi in elkaar passen of aan elkaar raken. Vaak moet er aangetoond worden dat bepaalde figuren even groot zijn of moet een bepaalde verhouding worden berekend. De meeste sangaku’s dateren van de ‘Edo’-periode (1600-1868), toen Japan bijna volledig van de rest van
de wereld afgesloten was. 817 sangaku’s zijn bewaard gebleven.

Wilden de wiskundigen de goden bedanken voor de inspiratie? Misschien. Maar waarschijnlijk was het ophangen van de meetkundige tabletten eerder een manier om resultaten te `publiceren’ en de bezoekers uit te dagen om de
problemen op te lossen. Veel Japanse wiskundigen reisden van tempel naar tempel om kennis te nemen van elkaars sangaku’s.

We geven hier een voorbeeld van een eenvoudige sangaku die je kunt oplossen met de stelling van Pythagoras. We zijn niet zeker of die echt in een tempel hangt, maar het is in elk geval in de stijl van de sangaku's. Je vindt nog veel meer sangaku's in Huvent (2008) en in van Lint en Breeman (2015).

Zoals steeds bij sangaku's, moet je de gegevens afleiden uit de figuur. De blauwe cirkel is ingeschreven in het midden van de groene halve cirkel. De rode cirkel raakt aan de middellijn, de groene halve cirkel en de blauwe cirkel.

\begin{kader}{Een sangaku}{}

    Wat is de verhouding van de stralen van de rode en de blauwe cirkel?

    \afbeelding[6cm]{img/LoepGeschiedenis_ges4.JPG}{Sangaku}{}
    
\end{kader}

Om het probleem op te lossen, geven we namen aan de middelpunten en aan enkele raakpunten (zie \autoref{ges5}). De straal van de blauwe cirkel noemen we $R$ en die van de rode cirkel $r$. De straal van de groene halve cirkel is dan $2R$. De stippellijnen van \autoref{ges5} drukken uit dat telkens als twee cirkels elkaar raken, het raakpunt en de twee middelpunten op één rechte liggen. Hierdoor weten we bv. dat $|AC|=2R-r$ (de straal van de groene halve cirkel min de straal van de kleine rode cirkel).

\afbeelding[6cm]{img/LoepGeschiedenis_ges5.JPG}{Dezelfde figuur met stippellijnen en benamingen}{ges5}

Als we de stelling van Pythagoras in verschillende driehoeken toepassen, vinden we:
\[
\begin{array}{lrl}
    \text{in }\Delta ADE: &(R+r)^2  &=|AD|^2+(R-r)^2 \\
    \Rightarrow &4Rr      &=|AD|^2\\
    \text{in }\Delta ADC: &(2R-r)^2 &=|AD|^2+r^2\\
    \Rightarrow &4R^2-4Rr &=|AD|^2\\
\end{array}
\]
Uit de tweede en vierde gelijkheid elimineren we \(|AD|\) en vinden we \(R = 2r\).
De gevraagde verhouding $\frac{r}{R}$ is dus $\frac{1}{2}$.

\clearpage
\section{Gauss en de regelmatige veelhoeken}
\subsection*{Waarom deze keuze?}
	De grote wiskundigen uit de klassieke oudheid krijgen de meeste belangstelling in de literatuur en bij uitbreiding ook in onze lessen: Euclides, Archimedes, ... De belangrijke kopstukken uit de achttiende en de negentiende eeuw krijgen die veel minder: Euler, Gauss, ... Het niet ter sprake brengen van Carl Friedrich Gauss (1777-1857) in de wiskundeles is van eenzelfde orde als het niet ter sprake brengen van Napoleon in de geschiedenisles. Beide tijdgenoten zijn elk in hun eigen werkdomein zeer invloedrijk geweest. 
	
	\subsection*{Bij welke gelegenheid?}
	In de derde graad komen er heel wat thema's aan bod die rechtstreeks in verband staan met Gauss. Zo worden stelsels van eerstegraadsvergelijkingen in het vijfde jaar opgelost door rijherleiden van een uitgebreide matrix via de methode van Gauss-Jordan. In de statistiek maken we gebruik van de normale verdeling en de kleinste-kwadratenmethode. Beide onderwerpen werden voor het eerst door Gauss beschreven. In de tweede en derde graad komen de complexe getallen aan bod. We stellen ze grafisch voor in het vlak van Gauss. En wanneer we het later hebben over de hoofdstelling van de algebra die zegt dat elke complexe veelterm van de $n$-de graad precies $n$ nulwaarden heeft, dan hebben we het eerste bewijs van deze stelling ook te danken aan Gauss.
	
	Bij elk van deze gelegenheden is het zinvol om het vernuft van Gauss in de kijker te zetten, ook al manifesteert dit vernuft zich niet altijd in onderzoeksdomeinen die tot de reguliere leerstof behoren. Tijdens het schrijven van deze loep lijkt de meetkunde van de cirkel uit de leerplannen geschrapt te zijn. We moeten echter zeker tijd kunnen maken voor een korte uitweiding over de uitzonderlijke prestatie van de 19-jarige Gauss die er in slaagde een constructie te bedenken om enkel met een passer en een latje een regelmatige zeventienhoek te tekenen, ingeschreven in een gegeven cirkel. 
	
	Een andere uitzonderlijke prestatie van Gauss is het terugvinden van het `grootste vermiste voorwerp sinds de oertijd'. Hij deed dit op 23-jarige leeftijd. Hierover verderop meer.
    \vfill\null
    \columnbreak
	
	\subsection*{Het leven van Gauss}
	
	Carl Friedrich Gauss werd in 1777 in een volks gezin geboren in Braunschweig, een stadje in het koninkrijk Hannover (nu centraal Duitsland). Zijn vader was tuinman en bouwopzichter, de moeder was een dienstmeid die geen onderwijs had genoten. Carl Friedrich zou wellicht ook geen hogere onderwijskansen hebben gekregen, mocht hij niet opgemerkt zijn door zijn mecenas, de hertog van Braunschweig, die hem een toelage gaf om op 15-jarige leeftijd naar het Collegium Carolinum van Braunschweig te gaan en op 18-jarige leeftijd naar de universiteit van Göttingen, een hondertal kilometer verderop.
	
	\afbeelding[7cm]{img/LoepGeschiedenis_gauss.jpg}{Carl Friedrich Gauss (1777-1855) op een bankbiljet van 10 Duitse Mark}{fig:gauss}
	
	Volgens een wijdverspreide anekdote verbaasde de jonge Carl Friedrich zijn meester in de lagere school al door in enkele ogenblikken een lange rekenopdracht uit te voeren waarmee de leerkracht de klas zoet probeerde te houden. Hij gaf hen de opdracht de getallen van 1 tot 100 op te tellen. Carl Friedrich had het inzicht om deze getallen twee aan twee op te tellen: de 1 met de 100, de 2 met de 99, \ldots en de 50 met de 51. Zo verkreeg hij in een oogwenk 50 sommen die gelijk waren aan 101. Volgens bepaalde biografieën toonde Gauss zijn antwoord (5050) aan de meester met de woorden `Ligget se' (hier heb je ze).
	
	
	Om onduidelijke redenen ruilde Gauss in 1798 de universiteit van Göttingen in voor de kleinere en minder gereputeerde universiteit van Helmstedt, dichter bij huis. Daar diende hij een jaar later zijn proefschrift in waarin hij de hoofdstelling van de algebra bewees over het aantal nulwaarden van een veeltermvergelijking van graad $n$.
	
	In 1801, op 23-jarige leeftijd, publiceerde hij zijn belangrijkste werk, de `disquisitiones arithmeticae' (Latijn voor `rekenkundig onderzoek'), waarin hij het modulorekenen introduceerde. Sommige onderwerpen hieruit had hij al jaren eerder uitgedacht. Zo kondigde hij al in 1796 aan een constructie gevonden te hebben voor de regelmatige zeventienhoek, maar gaf hij zijn geheimen pas prijs in de `disquisitiones arithmeticae'. Met deze vroegtijdige aankondiging veroverde Gauss als jonge snaak de bewondering van de brede wiskundige gemeenschap.
	
	Gauss heeft niet zo lang van zijn mecenas kunnen genieten als Simon Stevin destijds. Toen de hertog van Braunschweig in 1806 na verwondingen in de Slag van Jena overleed en Westfalen in handen viel van Jérôme Bonaparte (broer van Napoleon), kreeg Gauss het financieel moeilijk. Om zijn gezin te onderhouden, aanvaardde hij het ambt van directeur van het sterrenkundig observatorium van Göttingen. 
	
		
	Toen hij kort na het overlijden van zijn eerste vrouw hertrouwde en een nieuw gezin stichtte, aanvaardde hij ook de opdracht om heel Hannover op te meten door middel van driehoeksmeting. De naburige meetpunten die zich op kerktorens en bergtoppen bevonden, kon hij wederzijds zichtbaar maken door zijn pas uitgevonden heliotroop, een toestel dat zonlicht kon opvangen en nauwkeurig in een andere richting kon weerkaatsen. Om zijn opdracht als landmeter uit te voeren, was Gauss soms maanden van huis. Hij vond elke dag dat hij met landmeting bezig was een verlies voor de wetenschap, maar kon dit verlies ook relativeren.

 \afbeelding[7cm]{img/LoepGeschiedenis_heliotroop.jpg}{De heliotroop van Gauss op een bankbiljet van 10 Duitse Mark}{fig:heliotroop}
	
	In de jaren '30 van de negentiende eeuw was de kennis van de aarde nog niet wat ze nu is. Men twijfelde nog of de vorm van de aarde een oblate (pilvormige) of een prolate (sigaarvormige) ellipsoïde was. Er waren discussies of de aarde twee of vier magnetische polen heeft. Om uitsluitsel te geven, werden er verre expedities gedaan en werden er observatoria voor aardmagnetisme opgericht. 
	
	In 1833 bouwde Gauss samen met zijn goede vriend Max Weber het eerste magnetische observatorium in Göttingen. Rond die tijd knutselden Gauss en Weber ook de eerste elektromagnetische telegraaf in elkaar. In zijn briefwisseling was Gauss lyrisch over deze uitvinding en zelfs visionair: 
	\begin{citaat}
		De keizer van Rusland zou hiermee zijn bevelen binnen een minuut zonder tussenstations kunnen overbrengen van Sint-Petersburg naar Odessa.
	\end{citaat}
	
	Naar het einde van zijn leven werd Gauss minder productief. Hij had een zwakke gezondheid, hij leed aan astma en had aanvallen van depressiviteit. Carl Friedrich Gauss stierf op 23 februari 1855 aan een hartaanval.
	
	\subsection*{Regelmatige veelhoeken}  
	
	Niet elke regelmatige, in een cirkel ingeschreven $n$-hoek is enkel met een passer en een latje te construeren (dus volgens de regels van de Griekse oudheid). Dit is wel mogelijk voor \(n\) gelijk aan \(3, 4, 5, 6, 8, 10, \ldots\) In de \emph{Almagest} van Ptolemaeus (rond 100 na Christus) kunnen we deze constructies al terugvinden. We laten die van de vijfhoek, de zeshoek en de tienhoek hier even zien.
	
	Gegeven is de cirkel $c$ met middelpunt $M$ en met twee onderling loodrechte middellijnen $[BC]$ en $[DE]$, zie \autoref{fig:tienhoek1}. Deel het lijnstuk $[MC]$ in twee door middel van een hulpcirkel en een hulplijnstuk. Het punt $A$ is het midden van $[MC]$.
	
	\afbeelding[5.7cm]{img/LoepGeschiedenis_tienhoek1.jpg}{Het middelpunt van een straal}{fig:tienhoek1}
	
	Vervolgens teken je een hulpcirkel met middelpunt $A$ en met straal $|AE|$. Deze cirkel snijdt het lijnstuk $[BM]$ in het punt $H$ dat tussen $B$ en $M$ ligt. De driehoek $EHM$ is de basisdriehoek voor het tekenen van de ingeschreven regelmatige 5-hoek, 6-hoek en 10-hoek. De schuine zijde van deze rechthoekige driehoek is de zijde van de regelmatige 5-hoek, de lange rechthoekszijde is de zijde van de regelmatige 6-hoek en de korte rechthoekszijde is de zijde van de regelmatige 10-hoek.
	
	\afbeelding[6.3cm]{img/LoepGeschiedenis_tienhoek2.jpg}{De basisdriehoek}{fig:tienhoek2}
	
	Pas de zijden van deze veelhoeken nu af op de cirkelomtrek. In \autoref{fig:tienhoek3} deden we dit voor de 5- en de 10-hoek. Met de 6-hoek erbij zou de tekening te overladen zijn. Onder andere in (Kindt 2015) vind je een bewijs van de constructie met de basisdriehoek voor de 5-, 6- en 10-hoek.
	
	\afbeelding[6.3cm]{img/LoepGeschiedenis_tienhoek3.jpg}{De ingeschreven 5- en 10-hoek}{fig:tienhoek3}
	
	Sinds de oudheid bleef de theorie van de construeerbare regelmatige veelhoeken onaangeroerd. Dit bleef zo tot de 19-jarige Gauss de wereld na anderhalf millenium plots verbaasde met de aankondiging van een constructie voor een regelmatige 17-hoek, die pas later gepubliceerd zou worden wanneer de theorie helemaal uitgewerkt was. Vijf jaar later (in 1801) publiceerde hij effectief zijn \emph{Disquisitiones Arithmetica} waarvan het laatste gedeelte een volledig uitgewerkte studie was van de construeerbaarheid van de regelmatige $n$-hoeken.
	
	\afbeelding[5.9cm]{img/LoepGeschiedenis_zeventienhoek picuki.jpg}{Constructie van zeventienhoek volgens Gauss}{fig:zeventienhoek}
	
	Naast de constructie van de regelmatige 17-hoek bewees Gauss welke regelmatige $n$-hoeken er nog meer construeerbaar waren. Hiervoor had hij de priemgetallen van Fermat nodig. Deze priemgetallen zijn van de vorm $F_n=2^{2^n}+1$ met $n$ een natuurlijk getal. Momenteel zijn er slechts vijf priemgetallen van Fermat bekend. Maar in de tijd van Euler was men er van overtuigd dat er meer Fermatpriemgetallen waren. Euler toonde met zijn extreme rekenvaardigheid aan dat $F_5$ deelbaar was door 641. Een zeer kort en bevattelijk bewijs van de deelbaarheid van $F_5$ door 641 vind je in (Kindt 2015).
	
	\begin{tabel}
		\begin{tabular}{C{1cm} C{1cm} C{1cm}C{1cm} C{1.5cm}}
			\hlinethick
			\tableheader{$F_0$} & \tableheader{$F_1$} & \tableheader{$F_2$} & \tableheader{$F_3$} &\tableheader{$F_4$} \\
			\hlinethick
			3 & 5 & 17 & 257 & 65 537\\
			\hlinethick
		\end{tabular}
		\caption{De vijf gekende priemgetallen van Fermat}\label{tab:veer}
	\end{tabel}

 Door de passer-en-liniaalconstructie van de bissectrice van een hoek, kan het aantal zijden van een construeerbare regelmatige veelhoek verdubbeld worden. Als een vijfhoek construeerbaar is, is ook een tienhoek construeerbaar en een twintighoek.
	
	Gauss bewees dat de construeerbare regelmatige $n$-hoeken met een \emph{oneven} aantal zijden precies deze zijn waarvoor $n$ een product is van een aantal \emph{verschillende} priemgetallen van Fermat. Zo is de regelmatige 51-hoek construeerbaar want $51=3\cdot 17$. En ook de regelmatige 255-hoek want $255=3 \cdot 5 \cdot 17$. Maar de regelmatige 45-hoek is dit niet want $45=3 \cdot 3 \cdot 5$. De priemfactoren zijn hier niet verschillend.
	
	 
	
	Met de Fermatpriemgetallen die nu gekend zijn, blijft het aantal construeerbare regelmatige veelhoeken met een \emph{oneven} aantal zijden zeer beperkt. Het zijn er 31, tel maar na.
    Het is echter niet geweten of er nog andere Fermatpriemgetallen zijn dan de vijf gekende.
	
	Tot slot nog een klein staartje aan deze uitweiding. Kun je met een passer en een latje een hoek van $3^\circ$ kunt tekenen? En een hoek van $1^\circ$? Ziehier het antwoord op de eerste vraag. Een regelmatige 120-hoek is construeerbaar want $120=2^3 \cdot 3 \cdot 5$. Een $3\cdot5$-hoek is immers construeerbaar volgens het criterium van Gauss en je mag altijd vermenigvuldigen met een macht van $2$. Dit betekent dat een hoek van $3^\circ$ construeerbaar is want deze hoek is de middelpuntshoek die een zijde van de regelmatige 120-hoek ondersteunt. Voor het antwoord op de tweede vraag redeneren we uit het ongerijmde. Als de hoek van $1^\circ$ construeerbaar zou zijn, zou dit betekenen dat de 360-hoek te tekenen is met een passer en een latje. Maar dat is hij niet want $360=2^3 \cdot 3^2 \cdot 5$. De oneven priemfactoren zijn niet allemaal verschillend. De hoek van $1^\circ$ is dus niet construeerbaar.
 
	\subsection*{Het grootste vermiste voorwerp}
	Voor 1800 kenden de wetenschappers zeven planeten. Er waren gezaghebbende Duitse filosofen zoals Hegel die met wijsgerige argumenten verkondigden dat er nooit meer dan zeven planeten zouden gevonden worden. Op 1 januari 1801 echter ontdekte de onbekende Italiaanse astronoom Giuseppe Piazzi een achtste planeet vanuit zijn sterrenwacht in Palermo op Sicilië. Hij doopte de planeet Ceres naar de Romeinse godin van de akkerbouw (de Griekse godin Demeter). Voor de astronomen kwam deze ontdekking niet als een volledige verrassing. Er waren immers onverklaarbare afwijkingen opgemeten in de ellipstische banen van de planeten Mars en Jupiter. Wij kennen Ceres nu als de grootste planetoïde in de gordel tussen Mars en Jupiter.
 
	Piazzi volgde de baan van Ceres van 1 januari tot 11 februari. Daarna bleek ze vermist te zijn. Ze volgde echter keurig haar baan in de `schaduw van de zon'. Piazzi bezat op dat moment maar informatie over een boogsegment van $9^\circ$ van de baan van Ceres, wat volgens de toenmalige wetenschappers onvoldoende was om de hele baan vast te leggen.
	
	Tijdens de laatste voorbereidingen voor de druk van zijn \emph{Disquisitiones Arithmeticae} nam Gaus, die zeer geïnteresseerd was door deze mysterieuze verdwijning, de tijd om de baan van Ceres te voorspellen. Zijn berekening week erg af van de berekeningen van zijn collega's astronomen. Maar Gauss had het bij het rechte eind. Op oudejaarsavond herontdekte een amateurastronoom de planetoïde Ceres. De reputatie van Gauss als een briljant wetenschapper kreeg in 1801 in heel Europa dus een dubbele boost, door de publicatie van zijn belangrijkste werk en door het terugvinden van het grootste vermiste voorwerp ooit. De Franse wiskundige Laplace noemde hem `een buitenaardse geest in een menselijk lichaam'. Enkele maanden nadien kwam er een tweede planetoïde in beeld, Pallas. Gauss die er op gebrand leek om de betrouwbaarheid van zijn methode te bewijzen, berekende de baan in slechts enkele uren.
 \end{multicols}
 \afbeelding[9.6cm]{img/gauss.jpg}{}{}{}
 \clearpage
 \begin{multicols}{2}

	\section{Florence Nightingale en de beschrijvende statistiek}
	
	\subsection*{Het leven van Florence Nightingale}
	Florence Nightingale (1820-1910) was een Britse verpleegster die grondig de manier heeft veranderd waarop de maatschappij de verzorging van zieken en gewonden bekeek. Ze pionierde niet alleen op het gebied van de moderne verpleging, ook binnen het domein van de beschrijvende statistiek heeft ze een belangrijke bijdrage geleverd. Ze bedacht het zogenaamde pooldiagram om data duidelijk te visualiseren zodat die meer toegankelijk en begrijpelijk werden. Op die manier kon ze politici overhalen om meer te investeren in gezondheidszorg. Ze gebruikte statistische gegevens om de relatie tussen patiëntensterfte en hygiëne in ziekenhuizen te bewijzen en om de efficiëntie van ziekenhuizen te evalueren. Haar werk had een enorme invloed op de ontwikkeling van de moderne epidemiologie, de verpleegkunde en de volksgezondheid.
	
	\afbeelding[0.8\linewidth]{img/LoepGeschiedenis_florence1.png}{Florence Nightingale (1820-1910)}{fig:florence1} 	
	
		
	Florence Nightingale kwam uit een rijke familie van grootgrondbezitters en al vroeg tijdens haar jeugd bleek dat ze veel aanleg had voor wiskunde en vreemde talen. Voor de opvoeding van vrouwen uit haar sociale klasse werd wiskunde echter als onbelangrijk gezien en haar ouders keurden haar interesse dan ook af. Nightingale moest het daarom vooral van zelfstudie hebben. Ze verslond boeken over wiskunde en statistiek. 
 
 Gaandeweg kreeg ze ook interesse voor ziekenhuisbeheer. Ze wilde verpleegster worden, maar dat stond voor iemand van haar stand niet hoog genoeg aangeschreven. Nightingale overwon deze vooroordelen, zette tegen de wil van haar ouders door en werd toch verpleegster. Haar werk was haar grote passie en haar gedrevenheid was enorm. Ze gebruikte haar talent en kennis om anderen te helpen. Dit alles maakt haar een bijzonder inspirerende figuur.
	
	\subsection*{Bij welke gelegenheid?}
	
	Om de leerlingen te motiveren voor beschrijvende statistiek kun je de verwezenlijkingen van Florence Nightingale kort bespreken bij de start van de lessenreeks. Statistiek komt in elke graad van het secundair onderwijs voor. Het verhaal van Nightingale toont het belang en de kracht van beschrijvende statistiek.
	
	\subsection*{Data grafisch voorstellen: kort historisch overzicht}
	
	Naar het einde van de achttiende eeuw wilden steeds meer staten info vergaren over hun inwoners omdat ze beseften dat hun macht niet enkel afhing van de grootte van hun staat, maar ook van bijvoorbeeld de productiviteit van de bevolking. Overheidsambtenaren gingen systematisch data verzamelen zoals bevolkingsaantallen, info over de gezondheid van de bevolking, het inkomen, de prijsevoluties $\ldots$ In Duitsland ontstond zo de discipline `staatshuishouding' of `Statistik'. Het woord `statistiek' is afkomstig van het Latijnse `statisticum collegium', wat `les over staatszaken' betekent.	
	
	Een belangrijke innovatie in de negentiende eeuw was het gebruik van grafieken om al die gegevens te visualiseren. Verschillende wiskundigen leverden een bijdrage aan deze ontwikkeling. 
 
 William Playfair (1759-1823, niet te verwarren met zijn tijdsgenoot John Playfair naar wie het Playfair axioma van Euclidische meetkunde is vernoemd) was één van de eersten die tabellen met cijfers verving door visuele representaties. Hij ontwierp verscheidene soorten grafieken zoals lijn-, staart- en staafdiagrammen. Een bekend voorbeeld met zowel een lijn- als een staafdiagram zie je in figuur \ref{fig:playfair1}.
    \vfill
    \clearpage
    
	\afbeelding[\linewidth]{img/LoepGeschiedenis_playfair1.png}{Prijs van tarwe (staafdiagram) en gemiddelde weeklonen (lijngrafiek) }{fig:playfair1}
	
	De grafiek in figuur \ref{fig:playfair2} (1801) toont de relatie tussen belastingopbrengst, bevolkingsgrootte en oppervlakte van verschillende Europese landen. Een voorbeeld van multivariate analyse dus. De oppervlakte van de cirkelschijven representeert de oppervlakte van de genoemde landen. Het verticale lijnstuk links van elke cirkel stelt de bevolkingsgrootte van het land (in miljoenen) voor; het lijnstuk rechts is de belastingsopbrengst in miljoenen pond sterling. De verbindingslijnen tussen deze verticale rechten zeggen iets over het verband tussen deze twee.
 
 De tweede cirkelschijf visualiseert de data voor het Turkse Rijk. Deze schijf is verdeeld in drie sectoren die de locatie weergeven: Europa, Afrika en Azië. Het is het oudste bekende voorbeeld van een taartdiagram.
	
	\afbeelding[\linewidth]{img/LoepGeschiedenis_playfair2.png}{Het eerste taartdiagram}{fig:playfair2} 
	
	Het valt op dat deze grafieken ook voor ons nu nog steeds heel overzichtelijk en duidelijk zijn.
	
	Histogrammen werden voor het eerst gebruikt door A. M. Guerry in Frankrijk in 1833. Ze werden echter pas gemeengoed nadat Karl Pearson ze in 1895 gebruikte in een artikel over evolutieleer. Figuur \ref{fig:pearson1} is een voorbeeld uit dat artikel, waaruit duidelijk blijkt dat er minder bloemen zijn met een groter aantal bloemblaadjes.
	\afbeelding[\linewidth]{img/LoepGeschiedenis_pearson1.png}{Histogram, Karl Pearson (1895)}{fig:pearson1}
	
	Rond 1875 bedacht Francis Galton het ogief, om cumulatieve frequenties voor te stellen. In figuur \ref{fig:galton1} toont Galton de trekkracht van 519 mannen van 23 tot 26 jaar, in Engelse ponden (lbs). Hij plaatste de kracht langs de verticale as en het percentiel langs de horizontale as (dat doen wij nu meestal omgekeerd). Je kunt uit de grafiek bijvoorbeeld aflezen dat mannen uit het 70e percentiel een gewicht kunnen trekken van ongeveer 80 lbs. De grafiek is wat wij nu de inverse cumulatieve verdelingsfunctie noemen.
	\afbeelding[\linewidth]{img/LoepGeschiedenis_galton1.png}{Een reproductie van het eerste ogief, Galton (1875)}{fig:galton1}
	
	\subsection*{The lady with the lamp}
	
	Door haar kennis van wiskunde en statistiek én doordat ze al op jonge leeftijd veel boeken las over verpleegkunde en ziekenhuisbeheer, was Florence Nightingale al heel snel een expert binnen deze domeinen. Dit bracht haar snel in leidinggevende functies. Dank zij haar gegoede afkomst had ze bovendien connecties met politici uit de hoogste regionen en dat gaf haar invloed op het beleid.
	
	Tijdens  de  Krimoorlog van 1855, waar de Engelsen samen met de Fransen tegen de Russen vochten, kwam Florence Nightingale als diensthoofd  terecht  in het Scutari-ziekenhuis in de Bosporus waar ze gewonde Britse soldaten verzorgde.  Wat  zij  daar aantrof, was verschrikkelijk. Aan alles was er gebrek: schoon water, verband en medicijnen$\ldots$ De gewonden lagen met honderden naast elkaar op de vuile grond. Onder deze barre omstandigheden  probeerde  Nightingale  het  lot  van  de soldaten te verlichten en de medische zorg te verbeteren. 's Nachts deed ze met een lantaarn de ronde door de zalen om patiënten te bezoeken en om te inspecteren of alle hygiënische maatregelen die ze als diensthoofd in haar ziekenhuis oplegde, werden opgevolgd. Onder de soldaten stond ze daarom al snel bekend als `The lady with the lamp'. De Krimoorlog was de eerste oorlog waarbij journalisten aanwezig waren. Die brachten verslag uit over de dame met de lamp en stilaan werd Nightingale een nationale held.

	\afbeelding[\linewidth]{img/LoepGeschiedenis_nightingale3.png}{`Crimean War: Florence Nightingale with her lamp at a patient's bedside', Henrietta Rae (1891)}{fig:nightingale3} 	
	
	Om de regering ervan te overtuigen dat de meeste soldaten niet stierven als gevolg van oorlogshandelingen maar wél door het gebrek aan goede ziekenzorg en hygiëne, begon ze dagelijks gegevens bij te houden en zocht ze een manier om die data visueel en overzichtelijk voor te stellen. Ze maakte daarbij dankbaar gebruik van de ideeën van haar tijdgenoot Adolphe Quetelet (1796‐1874), Belgische statisticus en de grondlegger van sociale statistiek. Nightingale bestudeerde onder andere zijn werk ‘Physique sociale’, nam er delen van over en schreef er commentaar bij. Ze ging echter verder dan Quetelet. Ze besefte dat ze ministers niet kon overtuigen van de noodzaak om meer te investeren in gezondheidszorg door hen enkel naar getallen te laten kijken. Ze werd de uitvindster van het pooldiagram en introduceerde zo het gebruik van statistiek in de gezondheidszorg. 
	
	In figuur \ref{fig:pooldiagram} zie je het pooldiagram dat Nightingale bedacht, `Diagram of the causes of mortality in the army in the East'. Dit diagram geeft de jaarlijkse sterftecijfers (per $1000$ en per maand) 
	\begin{itemize}
		\item die optraden als gevolg van te voorkomen ziekten (in blauw), 
		\item die het gevolg waren van oorlogsverwondingen (in rood),
		\item die te wijten waren aan andere oorzaken (in zwart).
	\end{itemize}


In de korte lesactiviteit hieronder, stellen we enkele vragen die je samen met leerlingen bij dit diagram kunt bespreken.

\end{multicols}


\afbeelding[\linewidth]{img/LoepGeschiedenis_nightingale2.jpg}{Pooldiagram van Florence Nightingale uit 1858}{fig:pooldiagram} 	


\begin{lesactiviteit}{Het pooldiagram van Florence Nightingale}

\begin{vraagenantwoord}	

\vraag{Welke info staat er op het diagram en hoe moet je het diagram lezen?}
\antwoord{Het aantal slachtoffers dat er in een bepaalde maand gevallen is, is recht evenredig met de
oppervlakte van de cirkelsector. Elke kleur stelt een bepaalde categorie voor.
}

\vraag{Overlappen de verschillende oppervlakten met elkaar?}
\antwoord{Alle oppervlakten moet je meten vanuit het middelpunt, het zijn elk cirkelsectoren. De verschillende categorieën overlappen elkaar. In het diagram van 1855 kun je dit zien aan de ene sector waarbij rood aan de buitenkant zit.
}
\vraag{Welke oppervlakte is de grootste?}
\antwoord{De oppervlakte die meestal uitsteekt: zeer duidelijk het aantal slachtoffers door gebrek aan zorg.
}
\vraag{In oktober 1854 en april 1855 valt de zwarte oppervlakte samen met de rode. Wat 
betekent dit?}
\antwoord{Er zijn evenveel slachtoffers door oorlogsverwondingen als door andere oorzaken.
}
\vraag{ In januari en februari 1856 vallen de blauwe oppervlaktes samen met de zwarte. Wat betekent dit?} 
\antwoord{Er zijn evenveel doden door gebrek aan zorg als door alle andere oorzaken. De rode
categorie van de oorlogsgewonden is weggevallen.}
\end{vraagenantwoord}

\end{lesactiviteit}

\begin{multicols}{2}

Na de Krimoorlog richtte Florence Nightingale in 1860 een verpleegstersschool op in Londen. Het was een van de eerste opleidingen op dit gebied en de invloedrijkste. Ze werkte mee aan tal van rapporten en schreef o.a. ook een handboek voor verpleegkundigen. Veel van Nightingales suggesties met betrekking tot ziekenhuisarchitectuur en gezondheidszorg zijn ondertussen gemeengoed geworden: afzonderlijke afdelingen om het besmettingsgevaar te verkleinen, het belang van licht en goede ventilatie, het belang van preventie (‘gezondenzorg is even belangrijk als ziekenzorg’) en haar opvatting dat ziekenhuizen elk jaar hun resultaten zouden moeten publiceren zodat die onderling statistisch vergeleken	en	verbeterd	kunnen	worden. Nightingale kan gerust één van de grote sociale hervormers van haar tijd worden genoemd.
	
	\section{Rationale en irrationale verhoudingen bij Euclides}
	
	\subsection*{Waarom deze keuze?}
	
	Voor de Oude Grieken waren de enige getallen de natuurlijke getallen verschillend van 0. Hoe konden zij, met dit primitief getalbegrip, toch zo'n straffe wiskunde ontwikkelen? Je zou je kunnen afvragen of hun omtrekken, oppervlakten en volumes in de meetkunde beperkt bleven tot wat mooi uitkwam op natuurlijke getallen. Zeker niet. Naast getallen hadden ze (meetkundige) `grootheden'. Op een vernuftige manier slaagden Eudoxos en Euclides erin om met rationale en irrationale verhoudingen van grootheden te werken, zonder deze verhoudingen als getallen te bekijken. Niet eenvoudig, maar zo mooi en creatief dat we dit aan sterkere leerlingen niet willen ontzeggen.
	
	\subsection*{Bij welke gelegenheid?}
	
	In de tweede graad komen voor het eerst irrationale getallen ter sprake, bv. met het bewijs dat $\sqrt{2}$ irrationaal is. Misschien is dit een goede gelegenheid om te vermelden dat het getalbegrip een hele evolutie heeft meegemaakt. Wat wij nu getallen noemen, zijn niet altijd getallen geweest. Omdat het begrip reëel getal niet in één keer in al zijn finesses door de leerlingen verworven wordt en omdat dit fragment nogal moeilijk en subtiel is, zou dit fragment ook in de derde graad een plaats kunnen krijgen, bij wijze van differentiatie voor sterkere leerlingen.
	
	\subsection*{Getallen en verhoudingen bij de Oude Grieken}
	
	De Oude Grieken kenden wel meetkundige grootheden zoals lengte, oppervlakte en volume, maar ze bekeken deze grootheden niet als getallen. Ze bestudeerden verhoudingen van gelijksoortige grootheden: verhoudingen van twee lengtes of van twee oppervlakten of van twee volumes. 
	
	Sommige van die verhoudingen komen overeen met een verhouding van twee `getallen'. Zo verhoudt het volume van een kegel zich tot het volume van een cilinder met zelfde grondvlak en hoogte, zoals de getallen 1 en 3 zich tot elkaar verhouden. Wij bekijken deze verhouding nu als één getal: $\frac{1}{3}$, maar voor de Grieken bleven het twee getallen.
	
	Andere verhoudingen van gelijksoortige grootheden zijn niet gelijk aan een verhouding van (natuurlijke) getallen, bv. de verhouding van de diagonaal van een vierkant tot de zijde van dat vierkant. Wij bekijken deze verhouding als het irrationale getal $\sqrt{2}$, maar voor de Grieken was deze verhouding niet exact uit te drukken met twee getallen.
	
	\subsection*{Rationale verhoudingen bij Euclides}
	
	%We gaan niet in op de voorafgaandelijke definities van Euclides. In boek VII definieert hij een getal als een hoeveelheid, samengesteld uit eenheden.
	
	Neem bv. twee lengtes $a$ en $b$ waarvan wij zouden zeggen dat $\frac{a}{b}=\frac{3}{4}$. Hoe konden de Grieken dit uitdrukken zonder het getal $\frac{3}{4}$?
	Er zijn twee manieren om dat te doen. Ofwel zeg je dat er een lijnstukje bestaat dat precies $3$ keer in $a$ past en $4$ keer in $b$. De lengten $a$  en $b$ zijn \emph{onderling meetbaar}; je kunt die exact meten met een zelfde lijnstukje. Een tweede manier is zeggen dat 4 keer $a$ even lang is als 3 keer $b$ (\autoref{ges6}).
	
	\afbeelding{img/LoepGeschiedenis_ges6.JPG}{Rationale verhouding}{ges6}
	
	Twee rationale verhoudingen kunnen soms aan elkaar gelijk zijn. Als $a$ zich tot $b$ verhoudt zoals $c$ tot $d$, dan vormen $a, b, c, d$ (in die volgorde) een \emph{evenredigheid}. Als bv. $4a=3b$ en $a, b, c, d$ een evenredigheid vormt, dan moet ook $4c=3d$. Euclides veralgemeent dit (zie kader).
	
	Merk op dat de vier lengtes (of andere grootheden) daarom niet noodzakelijk alle vier onderling meetbaar moeten zijn. Bv. lijnstukken van lengtes $3$, $4$, $3\sqrt{2}$ en $4\sqrt{2}$ (zoals wij nu zouden zeggen) is een voorbeeld van een evenredigheid bestaande uit twee gelijke rationale verhoudingen, maar $3$ en $3\sqrt{2}$ zijn onderling onmeetbaar. De vier grootheden $a, b, c, d$ van een evenredigheid hoeven ook niet alle vier van dezelfde soort te zijn. Maar het is wel nodig dat de eerste twee en de laatste twee van dezelfde soort zijn. Als $a$ en $b$ niet van dezelfde soort zouden zijn (bv. een lengte en een oppervlakte), dan kun je er geen ongelijkheid zoals $ma<nb$ mee maken. Je kunt geen appels met citroenen vergelijken...
	
	\begin{kader}{Elementen Vii, 20}{Evenredigheid}
		
		Vier grootheden $a, b, c, d$, waarbij $a$ en $b$ onderling meetbaar zijn en $c$ en $d$ ook, vormen een evenredigheid wanneer
		
		$\forall m, n \in \mathbb{N}_0: ma=nb \Leftrightarrow mc=nd$
		
	\end{kader}
	
	Wij zouden zeggen: beide verhoudingen $\frac{a}{b}$ en $\frac{c}{d}$ zijn gelijk aan dezelfde breuk $\frac{n}{m}$.
	
	Merk op: de notaties in het kader, met kwantoren, de verzameling $\mathbb{N}_0$, de dubbele pijl... zijn niet die van Euclides, maar dit verandert niets aan de betekenis.
	
	\subsection*{Irrationale verhoudingen bij Euclides}
	
	Hoe definieerden de Oude Grieken een irrationale verhouding van gelijksoortige grootheden? Of om het op zijn Grieks te zeggen, voor twee grootheden die \emph{onderling onmeetbaar} zijn? Denk bv. aan lengtes $a$ en $b$ met $\frac{a}{b}=\sqrt{2}$, zoals dat het geval is als $a$  de diagonaal en $b$ de zijde is van een zelfde vierkant.
	
	Omdat er geen natuurlijke getallen $m, n$ bestaan zodat $\frac{a}{b}=\frac{n}{m}$ (zoals wij dit noteren), kunnen de Oude Grieken niet zeggen dat een veelvoud van $a$ gelijk is aan een veelvoud van $b$ zoals bij een rationale verhouding. Maar net zoals wij $\sqrt{2}$ door breuken kunnen \emph{benaderen}, kunnen zij ook werken met veelvouden $ma$ en $nb$ die bij benadering aan elkaar gelijk zijn. 
	
	$\sqrt{2}$ is ruwweg te benaderen door $\frac{3}{2}=1,5$. Een betere benadering is $\frac{7}{5}=1,4$, of nog beter $\frac{17}{12}\approx 1,4167$. Het is mogelijk om breuken te vinden die steeds dichter bij $\sqrt{2}\approx 1,41421$ liggen zonder er volledig mee samen te vallen. Gevorderde lezers herkennen hierin het feit dat $\mathbb{Q}$ \emph{dicht} is in $\mathbb{R}$.
	
	Omdat $\frac{3}{2}$ groter is dan $\sqrt{2}$, is $2$ keer $a$ kleiner dan $3$ keer $b$ (zie \autoref{ges70}). Omdat $\frac{7}{5}$ kleiner is dan $\sqrt{2}$, is $5$ keer $a$ groter dan $7$ keer $b$
	(zie \autoref{ges7}). Omdat $\frac{17}{12}$ groter is dan $\sqrt{2}$, is $12$ keer $a$ kleiner dan $17$ keer $b$. (Dit hebben we niet meer getekend omdat de benadering zo goed is dat je het verschil niet goed kunt zien.)
	
	\afbeelding[4cm]{img/LoepGeschiedenis_ges70.JPG}{Irrationale verhouding: ruwe benadering}{ges70}
 
	\afbeelding{img/LoepGeschiedenis_ges7.JPG}{Irrationale verhouding: betere benadering}{ges7}
	
	Euclides definieert in het algemeen een evenredigheid van vier grootheden, een gelijkheid tussen twee verhoudingen, ook wanneer het verhoudingen zijn van `onderling onmeetbare' grootheden zijn. (Denk bv. aan de diagonaal en de zijde van een vierkant en de diagonaal en de zijde van een groter of kleiner vierkant.)

 \begin{kader}{Elementen V,5}{Evenredigheid}
		Vier grootheden $a,\; b,\; c,\; d$, met $a,\; b$ van eenzelfde soort en $b,\;c$ van eenzelfde soort, vormen een evenredigheid wanneer
		$$\forall m, n\in\mathbb{N}_0: m a < n b \Leftrightarrow m c < n d$$
		$$\text{en\;} \forall m, n\in\mathbb{N}_0: m a > n b \Leftrightarrow m c > n d$$
	\end{kader}
	
	Hij zegt dat de verhoudingen van $a$ tot $b$ en van $c$ tot $d$ gelijk zijn, anders gezegd dat $a, b, c, d$ een evenredigheid vormen, wanneer voor alle mogelijke benaderingen geldt: een benadering die te klein is voor de ene verhouding, is ook te klein is voor de andere verhouding, en een benadering die te groot is voor de ene verhouding, is ook te groot is voor de andere verhouding.
	
	
	Opnieuw: Euclides noteerde het niet met deze symbolen, maar de essentie is dezelfde: hij definieerde een evenredigheid van willekeurige grootheden zonder andere dan natuurlijke getallen te gebruiken. Knap toch, hoe Euclides louter met natuurlijke getallen en ongelijkheden de irrationaliteit kon benaderen en `vangen'. Ere aan ere toekomt: deze theorie heeft Euclides waarschijnlijk overgenomen van zijn voorloper Eudoxox van Knidos (4de eeuw v.C.). Over het leven van Eudoxos is even weinig geweten als over dat van Euclides zelf. Er zijn geen rechtstreekse werken van Eudoxos bewaard gebleven.
	
	De definitie uit het tweede kader komt uit boek V van de elementen en die van het eerste kader uit boek VII. Vreemde volgorde bij Euclides? Dat komt omdat de boeken over meetkunde eerst komen en de studie van getallen pas begint in boek VII.
	
	Evenredigheden spelen een belangrijke rol in de meetkunde: denk maar aan gelijkvormige driehoeken, de stelling van Thales... maar ook oppervlaktes en volumes. In de zin ``De oppervlakten van twee driehoeken met een zelfde basis verhouden zich tot elkaar zoals de hoogtes'', heb je een evenredigheid $a, b, c, d$ met $a, b$ twee oppervlakten en $c, d$ twee hoogtes (dus lengtes). 
	
	\subsection*{Opmerking voor gevorderde lezers}
	
	Hoe Euclides irrationale verhoudingen aanpakte, komt voor een groot stuk overeen met de \emph{sneden van Dedekind} waarmee Richard Dedekind in de 19de eeuw de reële getallen `construeerde' vertrekkend van de rationale getallen. Het verschil is dat bij Dedekind rationale getallen al volwaardige getallen waren. Om bijvoorbeeld het irrationale getal $\sqrt{2}$ te `maken' verdeelde Dedekind de verzameling $\mathbb{Q}$ in twee deelverzamelingen: enerzijds de verzameling van de negatieve getallen en de positieve getallen waarvan het kwadraat kleiner is dan 2, en anderzijds de verzameling van de positieve getallen waarvan het kwadraat groter is dan 2. Het nieuwe getal $\sqrt{2}$ werd gedefinieerd als deze `snede' in de verzameling van de rationale getallen. Net als bij Euclides zijn twee nieuw-geconstrueerde irrationale getallen `gelijk' als dezelfde rationale getallen te klein zijn en dezelfde rationale getallen te groot...
	
	\section{Een kubus verdubbelen met passer, liniaal en parabolen}
	
	\subsection*{Waarom deze keuze?}
	
	Constructies met passer en liniaal zijn belangrijk in de Griekse meetkunde. Er is een wijdverspreide misvatting dat die constructies nodig zijn om nauwkeurige tekeningen te maken, want als je lengtes of hoeken meet, maak je meetfouten. In de praktijk maak je bij het werken met passer en liniaal ook fouten. Eerder dan over nauwkeurig tekenen, gaat het over de vraag wat je wel en wat je niet kunt verkrijgen door enkel rechten en cirkels met elkaar te snijden. Dat is het theoretische `spel' van de Griekse meetkunde. Om dit te begrijpen, volstaat het niet om enkel voorbeelden van constructies met passer en liniaal te bestuderen, zoals de constructie van de bissectrice van een hoek, maar is het essentieel om ook aandacht te hebben voor de tegenvoorbeelden: wat kun je met passer en liniaal \emph{niet} construeren? De drie klassieke constructie problemen zijn hierbij cultureel interessant: het heeft millennia geduurd voor men kon bewijzen dat er geen constructie mogelijk is. En bovendien zijn er wiskundigen die `out of the box of the ruler and compasses' gedacht hebben om deze problemen toch opgelost te krijgen, met andere hulpmiddelen dan passer en liniaal...
	
	\subsection*{Bij welke gelegenheid?}
	
	Deze paragraaf houdt verband met een derdegraadsvergelijking. Het zou dus eventueel passen bij veeltermfuncties in de derde graad. Het kan ook bij de studie van kegelsneden, bij de analytische meetkunde in de vrije ruimte.
	
	\subsection*{De verdubbeling van de kubus}
	
	Vooraf: we gaan ervan uit dat de leerlingen vertrouwd zijn met enkele basisconstructies met passer en liniaal, zoals de constructie van de loodlijn door een punt op een gegeven rechte, die van de evenwijdige door een punt aan een gegeven rechte.
	\clearpage
	
	\begin{kader}{Legende}{Verdubbeling van de kubus}
		
		Dit probleem waar de Oude Grieken hun tanden op braken, kreeg legendarische proporties. In de vijfde eeuw v.C. heerste een pestepidemie en er was uiteraard geen vaccin. Volgens de legende consulteerden de Grieken het orakel van Apollo op het eiland Delos om te weten hoe ze van de pest af konden geraken. Het orakel antwoordde dat ze hun kubusvormig altaar moesten verdubbelen. De ribbe verdubbelen was geen optie, want dan wordt het altaar (in volume) acht keer zo groot; het volume moest verdubbeld worden. Dat betekent dat de ribben moet worden vermenigvuldigd met $\sqrt[3]{2}$. Het probleem van de verdubbeling van de kubus wordt ook het \emph{Delische probleem} genoemd, naar het eiland Delos van het orakel.
		
	\end{kader}
	
	\afbeelding[5cm]{img/LoepGeschiedenis_Lverdubbeling.png}{Verdubbeling van de kubus}{}
	
	
	De derdegraadsvergelijking $x^3=2$ is voor ons geen probleem: de oplossing is $\sqrt[3]{2}$ en de kous is af. Maar voor de Oude Grieken was het niet zo eenvoudig. Ze bekeken deze vergelijking als het zoeken naar de ribbe van een kubus waarvan het volume het dubbel is van dat van een eenheidskubus. De oplossing van dit probleem kon geen getal zijn, want de getallen waren voor hen de natuurlijke getallen zonder nul: 1, 2, 3, 4, ... Ze zochten dus niet de `waarde' van die ribbe, maar een meetkundige constructie met passer en liniaal om die ribbe als lijnstuk te `maken' vertrekkend van een eenheidslijnstuk. 
	
	
	Een lijnstuk van lengte $\sqrt{2}$ (voor de Grieken: de zijde van een vierkant waarvan de oppervlakte het dubbel is van die van een eenheidsvierkant) kan wel met passer en liniaal geconstrueerd worden vanaf een eenheidslijnstuk. Het volstaat om met het eenheidslijnstuk als zijde, een vierkant te construeren (met de constructie van een loodlijn) en de diagonaal van dit vierkant te nemen. Ook een lijnstuk van lengte $\sqrt[4]{2}$ is te construeren vanaf een eenheidslijnstuk (zie vraag 5 in de lesactiviteit hieronder). ``Als we $\sqrt{2}$ en $\sqrt[4]{2}$ kunnen construeren, moet dit toch ook voor $\sqrt[3]{2}$ lukken", dachten de Grieken. Toch konden ze maar geen methode vinden om dit te doen. Veel later, in de 19de eeuw, heeft de Franse wiskundige Pierre Wantzel bewezen dat het met passer en liniaal alleen onmogelijk is.
	
	In de loop van de eeuwen is men vruchteloos blijven zoeken naar een constructie met passer en liniaal, maar sommige wiskundigen probeerden het probleem op te lossen door naast passer en liniaal, andere hulpmiddelen toe te laten. En dat lukte soms wel. Menaichmos, een leerling van Plato (4de eeuw v.C.) vond een constructie met passer, liniaal en twee parabolen (zie lesactiviteit). Diocles (rond 200 v.C.) gebruikte daarvoor een kromme, de cissoïde van Diocles (zie Roelens, Van den Broeck, 2022). Archytas (4de eeuw v.C.) vond een ingewikkelde oplossing door omwentelingslichamen in de ruimte met elkaar te snijden. In de twintigste eeuw werd ook het vouwen van papier (origami) ingezet als alternatieve manier om meetkundige constructieproblemen op te lossen. 
	
\end{multicols}

\begin{lesactiviteit}{Verdubbeling van de kubus met passer, liniaal en parabolen}

Een bekend probleem uit de Griekse Oudheid dat wel oplosbaar is met passer en liniaal, is de constructie van de middelevenredige van twee lengtes: gegeven twee lengtes $a$ en $b$, construeer een lengte $x$ zodat de verhouding van $a$ tot $x$ gelijk is aan de verhouding van $x$ tot $b$. Met modernere notaties:
$$\frac{a}{x}=\frac{x}{b}.$$

De oplossing gaat als volgt: je plaatst lijnstukken van lengte $a$ en $b$ achter elkaar. Je maakt een cirkel met diameter $a+b$. In het scheidingspunt $S$ tussen $a$ en $b$ teken je de loodlijn. Die snijdt de cirkel in $T$. Dan is $|ST|$ de gezochte middelevenredige.

\afbeelding[5cm]{img/LoepGeschiedenis_ges9.JPG}{}{}{}

\begin{vraagenantwoord}
	\vraag{Leg uit hoe deze constructie met passer en liniaal kan worden uitgevoerd.}
	\antwoord{Lijnstukken met lengtes $a$ en $b$ zijn gegeven. Eén van de twee verleng je en hierop pas je de lengte van de andere af, zodat ze achter elkaar staan. Het midden van het lijnstuk met lengte $a+b$ kun je construeren met de constructie van de middelloodlijn. In $S$ construeer je de loodlijn op het lijnstuk; daar kennen we ook een basisconstructie voor.}
	\vraag{Bewijs dat $|ST|$ inderdaad de middelevenredige is van $a$ en $b$.}
	\antwoord{De driehoeken $\Delta AST$ en $\Delta TSB$ zijn gelijkvormig (hoek-hoek).}
	
	\afbeelding[6cm]{img/LoepGeschiedenis_ges11.JPG}{}{}{}
	
	\vraag{Los die evenredigheid op naar $x$.}
	\antwoord{Dit geeft: $x=\sqrt{ab}$. Men noemt $x$ ook het meetkundig gemiddelde van $a$ en $b$. Het is een soort multiplicatieve versie van het rekenkundig gemiddelde: in plaats van op te tellen is het vermenigvuldigen en in plaats van te delen door twee is het de vierkantswortel nemen.}
	
	\vraag{Met de constructie van de middelevenredige kun je een lijnstuk construeren met lengte $\sqrt{a}$ als een eenheidslijnstuk en een lijnstuk van lengte $a$ gegeven zijn.}
	\antwoord{Neem $b=1$ in het antwoord van vraag 3. }
	
	\vraag{Leg uit: vertrekkend van een eenheidslijnstuk is een lijnstuk met lengte $\sqrt[4]{2}$ construeerbaar.}
	\antwoord{Je construeert eerst $\sqrt{2}$ en dan de middelevenredige van $1$ en $\sqrt{2}$.}
	
\end{vraagenantwoord}

Menaichmos (4de eeuw v.C.) bestudeerde een gelijkaardig probleem. Gegeven twee lengtes $a$ en $b$, zocht hij nu naar twee lengtes $x$ en $y$ zodat

$$\frac{a}{x}=\frac{x}{y}=\frac{y}{b}$$

Hij zocht met andere woorden naar de middelste twee termen van een meetkundig rijtje $(a, x, y, b)$ met $a$ en $b$ gegeven. We zouden $x$ en $y$ de twee \emph{tussenevenredigen} van $a$ en $b$ kunnen noemen. 

\begin{vraagenantwoord}[resume]
	
	\vraag{Toon aan: de tussenevenredigen $x$ en $y$ die Menaichmos zoekt, zijn de coördinaatgetallen van een snijpunt van de parabolen $x^2=ay$ en $y^2=bx$.}
	
	\antwoord{Met kruisproducten vinden we uit de eerste evenredigheid: $x^2=ay$ en uit de tweede: $y^2=bx$.}
	
	\vraag{Maak een schets van deze parabolen en bereken hiermee $x$ en $y$ in functie van $a$ en $b$.}
	
	\antwoord{We kwadrateren beide leden van de eerste gelijkheid: $x^4=a^2y^2$. We vervangen $y^2$ door $bx$ (tweede gelijkheid) en dat geeft: $x^4=a^2bx$. Dus: $x=0$ of $x^3=a^2b$. Op een analoge manier vinden we $y=0$ of $y^3=ab^2$. De twee snijpunten zijn $(0, 0)$ en $(\sqrt[3]{a^2b}, \sqrt[3]{ab^2)}$.}
	
	\afbeelding[8cm]{img/LoepGeschiedenis_ges10.PNG}{}{}
	
	\vraag{Leg uit hoe hiermee de verdubbeling van de kubus bereikt is.}
	\antwoord{Stel $a$ de ribbe van het oorspronkelijke altaar. We moeten een ribbe van lengte $x=a\sqrt[3]{2}$ construeren, de ribbe van een kubus met dubbel volume $2a^3$. Het volstaat twee parabolen te nemen zoals hierboven, maar met $b=2a$. Dan is de $x$-coördinaat van het snijpunt gelijk aan $\sqrt[3]{a^2 b}=\sqrt[3]{2a^3}=a\sqrt[3]{2}$.}
\end{vraagenantwoord}

Menaichmos zou een mechanisch toestel hebben gehad waarmee hij deze parabolen kon construeren. Met dit toestel samen met passer en liniaal is de verdubbeling van de kubus dus realiseerbaar.


\end{lesactiviteit}

\begin{multicols}{2}

\subsection*{De drie klassieke constructieproblemen}

Naast de verdubbeling van de kubus zijn er nog twee gelijkaardige constructieproblemen uit de Oudheid, die ook onoplosbaar bleken te zijn met passer en liniaal alleen.

\begin{itemize}
	\item De kwadratuur van een cirkel
	
	De Oude Grieken konden elke veelhoek met passer en liniaal omzetten in een vierkant met dezelfde oppervlakte, maar voor een cirkel lukt dit niet.
	
	\afbeelding[5cm]{img/LoepGeschiedenis_Lkwadratuur.png}{Kwadratuur van de cirkel}{}
	
	``Gegeven een cirkel, construeer een vierkant met dezelfde oppervlakte als die cirkel.'' 
	
	\item De trisectie van een hoek
	
	Terwijl een willekeurige hoek met passer en liniaal gemakkelijk in twee of in vier gelijke delen kan worden verdeeld, lukt het niet om die in drie gelijke delen te verdelen.
	
	\afbeelding[5cm]{img/LoepGeschiedenis_Ltrisectie.png}{Trisectie van een hoek}{}
	
	``Gegeven een hoek, verdeel die hoek in drie gelijke hoeken.''
	
	
\end{itemize}

\subsection*{Construeerbare getallen}

Als je een bepaalde lengte kunt construeren vanaf de eenheid, dan zeggen we dat het maatgetal van deze lengte een \emph{construeerbaar getal} is. Er zijn construeerbare en niet-construeerbare reële getallen. 

Descartes (17de eeuw) toonde aan dat een (reëel) getal construeerbaar is als en slechts als het getal verkregen kan worden door gehele getallen op te tellen, af te trekken, te vermenigvuldigen, te delen en te vierkantswortelen. Zo is $\sqrt[4]{1+\sqrt{2}}=\sqrt{\sqrt{1+\sqrt{2}}}$ construeerbaar en verrassend genoeg ook $\sqrt[3]{7+5\sqrt{2}}$ want $\sqrt[3]{7+5\sqrt{2}}=1+\sqrt{2}$ (zoals je kunt narekenen door van beide leden de derdemacht te nemen). Net zoals Menaichmos voerde hij ook extra krommen toe aan de cirkels en de rechten om de klassieke Griekse constructieproblemen met extra hulpmiddelen op te lossen. Deze krommen maakte hij met mechanische toestellen met stokken en vijzen, denk aan het spel Meccano. Maar hij beschreef die krommen ook met vergelijkingen; dit vormde het begin van de analytische meetkunde.	
	
	\section{De paradox van Bertrand en het kansbegrip in de twintigste eeuw}

	\subsection*{Waarom deze keuze?}
	
	Gedurende lange tijd werd kansrekening gedaan op basis van experimenten met een eindig aantal uitkomsten. Ook situaties die niet direct op die manier konden beschreven worden, zoals bijvoorbeeld experimenten met oneindig veel uitkomsten, paste men in dat kader (gebaseerd op het werk van Laplace). Dit leidde gaandeweg tot een aantal paradoxen en andere moeilijkheden in de theorie. Die hadden ook te maken met verschillende interpretaties van het kansbegrip. Er rezen serieuze twijfels bij de fundamenten van de kansrekening en de toepasbaarheid ervan in de `echte wereld'. Er was nood aan een strikt kader. 
	
	In deze paragraaf hebben we het over misschien wel de bekendste van de paradoxen binnen kansrekening: de paradox van Bertrand. We bekijken ook een verwant meetkundig kansprobleem. Paradoxen hebben niet alleen een rol gespeeld in de ontwikkeling van het kansbegrip, ze zijn ook gewoon verrassend en intrigerend. Het is bovendien zinvol dat leerlingen beseffen dat er tot op de dag van vandaag discussies zijn over het kansbegrip. Kansberekening (en bij uitbreiding wiskunde) is mensenwerk, een spel van conventies en afspraken. Afhankelijk van de afspraken krijg je soms andere resultaten.

	Vooraleer we de paradox bestuderen, beschrijven we kort nog de geschiedenis van  kansrekening. Het ontstaan van de moderne kansrekening situeert men vaak bij gokspelen. De Italiaanse arts en wiskundige Girolamo Cardano (1501-1576), die als eerste algemene oplossingen van de derde- en vierdegraadsvergelijking publiceerde, was een fervente liefhebber van het dobbelspel. Hij schreef een werk waarin hij enkele kansvragen behandelt o.a. bij het dobbelen met meerdere dobbelstenen. Een eeuw later, in 1654, legt de Franse ridder De Méré zijn ondertussen bekende partijenprobleem voor aan de wiskundige Blaise Pascal. Stel dat bij een dobbelspel tussen twee personen de speler die het eerst bijvoorbeeld zes partijen wint, de pot krijgt. Hoe moet de pot dan verdeeld worden wanneer een dobbelspel tussen deze twee spelers voortijdig wordt afgebroken? Je vindt dit probleem vandaag terug in verschillende handboekenreeksen. 
	
	Blaise Pascal correspondeerde hierover met Pierre De Fermat en uit die correspondentie ontstond een heel nieuw vakgebied binnen wiskunde. In 1812 verscheen het werk \textit{Théorie analytique des probabilités} van de Franse wiskundige Pierre-Simon Laplace. Hij geeft daarin een meer rigoureuze definitie van het kansbegrip en hij beschrijft de kanstheorie zoals die tot dan ontwikkeld is.
	
	Kansrekening bleek niet alleen toepasbaar bij gokspelen, er volgen ook succesvolle toepassingen binnen astronomie (o.a. om grip te krijgen op meetfouten bij observaties van hemellichamen) en statistische fysica (bij het beschrijven van eigenschappen van gassen aan de hand van random bewegingen van grote aantallen deeltjes). Vanaf het begin van de twintigste eeuw worden de toepassingsdomeinen breder dan de exacte wetenschappen en gebruikt men kansrekening bij experimenten en modellen in economie, geneeskunde, biologie, psychologie$\ldots$ 
		
	In 1933 publiceerde de Russische wiskundige Andrej Kolmogorov (1903-1987) het boek \textit{Grundbegriffe der Wahrscheinlichkeitsrechnung} waarin hij een axiomatische aanpak van de kansrekening voorstelt, daarbij steunend op verzamelingenleer en maattheorie. Hij legde hiermee de basis voor de kansrekening zoals wij die nu kennen. De rigoureuze aanpak van Kolmogorov bood een antwoord op de vragen rond de paradoxen en de fundamenten van de kansrekening. Dat neemt niet weg dat er tot op vandaag (filosofische) discussies zijn over de paradox van Bertrand (zie verderop) en de interpretatie van het kansbegrip.

 Interessant weetje voor de leraar: op het hoogtepunt van zijn wetenschappelijke carrière koos Kolmogorov ervoor om zich volledig en met veel overtuiging op zijn lesopdracht te storten. Hij stopte met wetenschappelijk onderzoek en het laatste kwart van zijn leven wijdde hij volledig aan lesgeven. 
    \afbeelding[0.95\linewidth]{img/LoepGeschiedenis_kolmogorov.png}{Andrej Kolmogorov (1903-1987)}{fig:Kolmogorov}
 
		
	\subsection*{Bij welke gelegenheid?}
	Dit onderdeel past bij de lessen kansrekening van de derde graad.
	
	In richtingen met weinig wiskunde of als er weinig tijd is, kun je een schets geven van de paradox en van de oplossing van het meetkundig kansprobleem zonder verder in te gaan op wiskundige details. In een richting met gevorderde wiskunde kun je een grondigere uitwerking tonen. Iets vertellen over deze paradox en de impact ervan op de ontwikkeling van het kansbegrip in de twintigste eeuw geeft ook de gelegenheid om het te hebben over axiomasystemen. In paragraaf \ref{euclides} hadden we het al over \textit{De Elementen} van Euclides en de axiomatische opbouw van de meetkunde die in dit werk wordt beschreven. Kolmogorov vond dat de kansrekening `op precies dezelfde manier zou moeten worden geaxiomatiseerd als de meetkunde en de algebra'.
	
	\subsection*{De paradox van Bertrand}
	De paradox is door de wiskundige Joseph Bertrand in 1889 voorgesteld in zijn werk \textit{Calcul des probabilités}.
	
	\begin{kader}{}{Paradox van Bertrand}
		Beschouw een cirkel en een gelijkzijdige driehoek ingeschreven in die cirkel. Als je willekeurig een koorde van de cirkel kiest, wat is dan de kans dat die koorde langer is dan de zijde van de driehoek?
	\end{kader}
	 
	 Bertrand kwam met drie verschillende en schijnbaar correcte oplossingsmethoden. De paradox bestaat erin dat elk van de methoden een ander antwoord geeft.
	
	\subsubsection*{Methode 1}
	We kiezen een willekeurige koorde door willekeurig twee punten van de cirkel te kiezen en die te verbinden.
	
	Zonder verlies van algemeenheid kunnen we één van de punten van de koorde laten samenvallen met een hoekpunt van de gelijkzijdige driehoek. Je kunt de cirkel immers altijd draaien zodat dat het geval is en dit verandert de lengtes van de koorde niet. Dit samenvallend punt is het rode punt op figuur \ref{fig:methode1}.	
	
	\afbeelding[0,7\linewidth]{img/LoepGeschiedenis_methode1.png}{Methode 1}{fig:methode1}
	
	De hoekpunten van de gelijkzijdige driehoek verdelen de cirkel in drie gelijke delen. Enkel wanneer je het tweede punt kiest op de cirkelboog die rood is gekleurd op figuur \ref{fig:methode1}, krijg je een koorde die langer is dan de zijde van de driehoek. De gevraagde kans is dus $\frac{1}{3}$.

 \vfill
 \clearpage
	
	\subsubsection*{Methode 2}
	We kiezen nu een koorde door eerst willekeurig de richting ervan te kiezen. 
	
	Draai nu de cirkel totdat één van de zijden van de driehoek dezelfde richting heeft en teken dan de straal die loodrecht staat op die zijde, zie figuur \ref{fig:methode2}. Merk op dat de zijde dan de middelloodlijn is van die straal.
	
	\afbeelding[0,7\linewidth]{img/LoepGeschiedenis_methode2.png}{Methode 2}{fig:methode2}
	
	Kies willekeurig een punt op de straal en teken de koorde met de gekozen richting die door dat punt gaat. Enkel wanneer het gekozen punt dichter ligt bij het middelpunt dan bij de rand (op het rode lijnstuk in figuur \ref{fig:methode2}), is de koorde langer dan de zijde van de driehoek. Op deze manier vinden we dat de gevraagde kans gelijk is aan $\frac{1}{2}$.
	
	 \subsubsection*{Methode 3}
	
		
	\afbeelding[0,7\linewidth]{img/LoepGeschiedenis_methode3.png}{Methode 3}{fig:methode3}

 \vspace{3cm}


 


	Een koorde ligt volledig vast als je het middelpunt ervan vast legt. We bepalen nu een willekeurige koorde door willekeurig een punt te kiezen binnen de cirkel. 

 	Teken dan de ingeschreven cirkel van de driehoek. De koorde die je bepaalde door het middelpunt te kiezen, is enkel langer dan de zijde van de driehoek als het gekozen middelpunt binnen de ingeschreven cirkel ligt (het rode gebied in figuur \ref{fig:methode3}).
	
	Stel dat de straal van de omgeschreven cirkel 1 is, dan kun je narekenen dat de straal van de ingeschreven cirkel $\frac{1}{2}$ is. De gevraagde kans is gelijk aan de verhouding van de oppervlakten van de twee cirkels en die verhouding is $\frac{1}{4}$.
	
	
	\subsubsection*{Conclusie}
	We moeten voorzichtig zijn wanneer het gaat om experimenten waarbij we een object willekeurig moeten kiezen. De oplossing hangt hier af van de manier waarop je de willekeurige koorde kiest. Elke methode geeft als het ware meer kans aan sommige koorden dan aan andere. 
	
	Het probleem van Bertrand is eigenlijk niet goed gesteld omdat de notie van `at random een koorde van een cirkel kiezen' niet nauwkeurig genoeg gedefinieerd is. Bij de formulering van het probleem zouden we moeten afspreken hoe het willekeurig kiezen van een koorde van een cirkel moet worden opgevat. Op het YouTubekanaal Numberphile vind je een heel duidelijk filmpje (Sanderson, 2021) met meer uitleg over de paradox van Bertrand.
	
	\subsubsection{Een spaghetto breken}
	Om kansproblemen op te lossen met experimenten met continue variabelen en oneindig veel uitkomsten, stelt men die uitkomsten soms voor als meetkundige objecten. Een maat voor het aantal uitkomsten is dan de lengte, oppervlakte, volume $\ldots$ van die objecten. We illustreren dit principe in de volgende lesactiviteit.

 \vfill
 \clearpage

	
\end{multicols}

\begin{lesactiviteit}{Kansproblemen meetkundig oplossen}
	In deze opdracht lossen we enkele kansproblemen op door uitkomsten meetkundig voor te stellen en het aantal uitkomsten te formuleren in termen van lengte, oppervlakte, volume$\ldots$\\
	
	\begin{kader}{}{Probleem 1}
		Wat is de kans dat een willekeurig reëel getal tussen $0$ en $3$ dichter bij $0$ ligt dan bij $1$?
	\end{kader}
	
	We stellen alle mogelijke uitkomsten, dit zijn alle reële getallen tussen $0$ en $3$, voor op een  getallenas:
	
	\afbeelding[0,45\linewidth]{img/LoepGeschiedenis_as1.png}{}{fig:as1}
	 	

	\begin{vraagenantwoord}	
		\vraag{Duid aan waar de gunstige uitkomsten liggen op de as.}
		\antwoord{Een getal ligt dichter bij $0$ dan bij $1$ als dat getal tussen $0$ en $0,5$ ligt.
		 \afbeelding[0,5\linewidth]{img/LoepGeschiedenis_as2.png}{}{fig:as2}}

		\vraag{Bereken nu de gevraagde kans.}
		\antwoord{We kunnen de lengte van de twee aangeduide lijnstukken gebruiken als maat voor het aantal gunstige en het aantal mogelijke uitkomsten. Met de regel van Laplace vinden we dat de gevraagde kans gelijk is aan de verhouding van die twee lengten: $\frac{0,5}{3}=\frac{1}{6}\approx 17\%$.}
	\end{vraagenantwoord}

	\vspace{0.5cm}
	\begin{kader}{}{Probleem 2}
		Zowel jij als de bus komen op willekeurige tijden tussen 12.00 uur en 13.00 uur aan bij de bushalte. Als de bus arriveert, wacht hij nog 5 minuten aan de halte voordat hij vertrekt. Als jij aankomt, wacht je 20 minuten voordat je vertrekt als de bus niet komt. Wat is de kans dat je de bus haalt?
	\end{kader}

	 Stel $B$ en $Y$ de tijden in minuten (vanaf 12.00 uur) totdat de bus, respectievelijk jij, arriveert aan de bushalte. Aangezien het hier gaat om twee onafhankelijke continue variabelen, formuleren we dit probleem als een tweedimensionaal meetkundig probleem. Alle mogelijke uitkomsten zijn punten binnen een vierkant met zijde 60 in het vlak:
	 
	\afbeelding[0,3\linewidth]{img/LoepGeschiedenis_vlak1.png}{}{fig:vlak1}
	  
	 
	 \begin{vraagenantwoord}	
		\vraag{Aangezien de bus 5 minuten wacht, moet jij arriveren binnen 5 minuten na de aankomst van de bus aan de bushalte, wil je de bus halen. Vertaal dit in termen van $B$ en $Y$.}
		\antwoord{We moeten eisen: $Y\leq B+5$.}
		
		\vraag{Welk gebied van het vierkant voldoet aan de voorwaarde die je vond bij de vorige vraag? Kleur dit gebied.}
		\antwoord{Het stuk onder de rechte $Y=B+5$ voldoet.\\ 	\afbeelding[0,32\linewidth]{img/LoepGeschiedenis_vlak2.png}{}{fig:vlak2}}
		
		\vraag{Jij wacht maar 20 minuten. Om de bus te halen, mag je dus ook niet meer dan 20 minuten voor de bus aankomen aan de halte. Vertaal dit in termen van $B$ en $Y$ en kleur het deel van het gebied dat voldoet.}
		\antwoord{De voorwaarde wordt nu $Y\geq B-20$. Het stuk boven de rechte $Y=B-20$ voldoet.\\
		\afbeelding[0,35\linewidth]{img/LoepGeschiedenis_vlak3.png}{}{fig:vlak3}}

		\vraag{Combineer de vorige twee antwoorden en duid alle gunstige punten (uitkomsten) aan binnen het vierkant van alle uitkomsten.}
		\antwoord{\afbeelding[0,33\linewidth]{img/LoepGeschiedenis_vlak4.png}{}{fig:vlak4}}
		\vspace{0.3cm}
		
		\vraag{Bereken nu de gevraagde kans.}
		\antwoord{We gebruiken de oppervlakte van het groene gebied als maat voor het `aantal' gunstige uitkomsten en de oppervlakte van het vierkant als maat voor het aantal mogelijke uitkomsten. De gevraagde kans is de verhouding tussen beide: $\displaystyle \frac{60^2-\frac{55^2}{2}-\frac{40^2}{2}}{60^2}=\frac{103}{288}\approx 36\%$.}		
	\end{vraagenantwoord}


	\vspace{0.5cm}		
	\begin{kader}{}{Probleem 3: de gebroken spaghetto}
		We laten een spaghetto vallen en die breekt op twee willekeurige plaatsen. Wat is de kans dat de drie stukken die je zo bekomt een driehoek vormen?
	\end{kader}

 

	
	\begin{vraagenantwoord}
	\vraag{We starten met een herhaling: ken je de driehoeksongelijkheid nog?}
	\antwoord{De som van twee zijden van een driehoek is strikt groter dan de derde zijde.}
	
	\vraag{Los het kansprobleem nu op.}
	\antwoord{Zonder verlies van algemeenheid kunnen we stellen dat de spaghetto lengte 1 heeft. Stel dat twee stukken van de gebroken spaghetto lengte $a$ en $b$ hebben. Het derde stuk heeft dan lengte $1-a-b$. Zonder verlies van algemeenheid stellen we: $a<b$. \\	
	\afbeelding[0,35\linewidth]{img/LoepGeschiedenis_spaghetto1.png}{}{fig:spaghetto1}  
	Dan geldt al zeker: $0<a<1$, $0<b<1$ en $0<1-a-b<1$. Deze ongelijkheden bakenen een gebied af in het vlak, de gekleurde driehoek in de figuur hieronder. De ongelijkheid  $0<1-a-b<1$ geeft hierbij aan dat het gebied begrensd is door de schuine zijde van de driehoek.
	\afbeelding[0,33\linewidth]{img/LoepGeschiedenis_spaghetto2.png}{}{fig:spaghetto2}
	Elk punt van dat gebied stelt een gebroken spaghetto voor, met een bepaalde waarde voor $a$, $b$ en $1-a-b$.\\ Welke punten komen overeen met gunstige uitkomsten (d.w.z. stukken die een driehoek vormen)? Opdat de stukken een driehoek vormen, moet voor elke zijde de driehoeksongelijkheid gelden:
	\hspace*{0.5cm} $a+b>1-a-b$ of nog $a+b>\frac{1}{2}$,\\
	\hspace*{0.5cm} $a+1-a-b>b$ d.w.z. $b<\frac{1}{2}$,\\ \hspace*{0.5cm} $b+1-a-b>a$ of nog $a<\frac{1}{2}$.\\
	We tekenen deze gunstige uitkomsten:
	\afbeelding[0,34\linewidth]{img/LoepGeschiedenis_spaghetto3.png}{}{fig:spaghetto3}
	Elk punt binnen de blauwe driehoek komt overeen met een gebroken spaghetto waarvan de stukken een driehoek vormen. De oppervlakte van deze blauwe driehoek is een maat voor het aantal gunstige uitkomsten. De oppervlakte van de grote gele driehoek is een maat voor het totaal aantal uitkomsten. De gevraagde kans is de verhouding van die oppervlakten en die is $\frac{1}{4}$.\\ De kans dat een spaghetto in drie willekeurige stukken breekt die een driehoek vormen, is dus $\frac{1}{4}$.}	
	\end{vraagenantwoord}	

  \afbeelding[14cm]{img/spaghetti.jpg}{}{}{}
	
\end{lesactiviteit}

	
\begin{multicols}{2}
	
We hebben bij de problemen uit de lesactiviteit hierboven de notie `at random' telkens nauwkeurig gedefinieerd. Bovendien zijn lengte, oppervlakte, volume$\ldots$ goed gedefinieerd. Zo komt het willekeurig breken van een spaghetto overeen met het kiezen van een punt op een driehoek en hebben we met de oppervlakte van zo'n driehoek een goede maat voor het aantal uitkomsten. Er zijn echter ook andere manieren van breken mogelijk. Je zou bijvoorbeeld het eerste punt kunnen kiezen op het lijnstuk, dan een muntstuk opgooien om te bepalen op welk van beide stukken je het tweede punt kiest, om daarna het tweede punt te kiezen. Dit levert andere kansen. 

Maattheorie geeft een rigoureus raamwerk voor kansrekening, inclusief kansen op oneindige verzamelingen. We gaan er hier niet verder op in.

Tot slot geven we nog een mooie andere meetkundige kansvraag mee.

\begin{kader}{}{Kans op een stompe driehoek}
Als je willekeurig drie punten kiest in het vlak, vormen ze (meestal) de hoekpunten van een driehoek. Wat is de kans dat deze driehoek stomphoekig is?
\end{kader}

Laat je leerlingen het antwoord eerst intuïtief bepalen of laat ze een gokje doen en laat ze pas daarna een oplossing beredeneren of berekenen. Een volledige uitwerking, mét voorstel voor een experimentje in de klas, vind je bij De Naeghel (2017).

Een analoge leuke toepassing waarbij kansen worden bepaald via oppervlakten, is de naaldproef van Buffon. Ook daar hoort een experimentje bij: je kunt de leerlingen $\pi$ laten benaderen door tandenstokers op een vel papier met lijnen te laten vallen. We gaan er hier niet verder op in, maar we maken er zeker ooit een Spinnenweb van.

	\section{Euler en de convergentie van reeksen}
	
	\subsection*{Waarom deze keuze?}
	
	Meer dan 300 jaar geleden werd de wellicht beroemdste Bazelaar (met een grote letter b, geen beuzelaar of raaskaller dus) en de wellicht ook de meest productieve wiskundige aller tijden geboren: Leonhard Euler. Zijn naam komt zo vaak terug in diverse domeinen van de wiskunde dat we ook in onze lessen niet aan hem kunnen voorbij gaan.
	
	\afbeelding[5cm]{img/LoepGeschiedenis_euler.jpg}{Leonhard Euler (1707-1783)}{fig:euler}
	
	In de meetkunde associëren we Euler met de rechte in een willekeurige driehoek die door de punten $H$ (het hoogtepunt), $M$ (het middelloodpunt) en $Z$ (het zwaartepunt) gaat. Euler stond aan de wieg van de grafentheorie. De Euler-graaf is een een graaf waarbij er een gesloten wandeling bestaat die alle paden precies één keer doorloopt. Ook verbinden we zijn naam aan de formule $Z-R+H=2$, die het verband aangeeft tussen het aantal zijvlakken, het aantal ribben en het aantal hoekpunten van een willekeurig veelvlak zonder tunnels (een lichaam met `geslacht' of `genus' 0). Elk veelvlak met één doorboring voldoet aan  $Z-R+H=0$.
	
	Ook was het Euler die de notatie van het \emph{exponentiële getal} e introduceerde. In verband hiermee: de identiteit van Euler, $e^{i\pi}+1=0$, die een verband legt tussen de vijf belangrijkste getallen e, $\pi$, i, 1 en 0, werd in verschillende polls in wiskundekringen aangewezen als mooiste formule uit de hele wiskundegeschiedenis.
	
	\subsection*{Bij welke gelegenheid?}
	
	In deze paragraaf bespreken we enkel de inbreng van Euler bij het onderzoek naar convergentie van reeksen (of sommen van rijen). Convergentie van reeksen komt onder andere ter sprake bij de berekening van de oppervlakte onder een kromme via  benaderingen met rechthoekjes (Riemannsommen). Vaak neem ik bij dit leerstofonderdeel de tijd om over deze markante wiskundige uit te weiden.
	
	\subsection*{Het leven van Euler}
	
	Euler werd in 1707 geboren in Reichen, een dorp in de buurt van Bazel, als zoon van een Calvinistisch predikant. Euler zelf was door zijn afkomst lotsbestemd predikant te worden. Evenals zijn vader volgde hij wiskundelessen bij een van de Bernoulli's, een gerenommeerde geslacht van wiskundigen uit Bazel. Toen Euler op 17-jarige leeftijd zijn `masterdegree' in de filosofie behaalde aan de universiteit van Bazel, pleitte leermeester Johann Bernoulli bij Eulers vader om zijn zoon verder te laten studeren in de wiskunde. Dat zou snel gebeuren, maar niet aan de universiteit van Bazel. Tijdens de Bazelse periode communiceerde Euler met de Bernoulli's onder andere over het onopgeloste probleem van de convergentie van de reeks $\frac{1}{1^2}+\frac{1}{2^2}+\frac{1}{3^2}+\frac{1}{4^2}+ \dots $, dat nu bekend staat als het \emph{probleem van Bazel}.
	
	Toen Euler 20 was, verliet hij Zwitserland. Hij diende de rest van zijn leven onder de drie Groten der aarde, achtereenvolgens: Peter de Grote (Sint-Petersburg), Frederik de Grote (Berlijn) en Catharina de Grote (Sint-Petersburg).
	
	Euler trok van Bazel naar de pas opgerichte universiteit van Sint-Petersburg. Dat deed hij nadat twee van de jongste telgen uit de Bernoulli-familie, Daniël en Nicolaus, ook al  geheadhunt waren door deze universiteit. Er was voor Euler enkel een leerstoel vrij in de faculteit van geneeskunde. Euler bereidde zich gedurende drie maanden voor op dit ambt en vertrok in 1727 naar Rusland. 
	
	Euler maakte snel naam als een briljant wiskundige. Hij ontkrachtte een vermoeden van Fermat, dat stelde dat elk natuurlijk getal van $F_n=2^{2^n}+1$ een priemgetal is. Euler toonde aan dat $F_5$, een getal van tien cijfers, deelbaar is door 641. Momenteel zijn er geen priemgetallen $F_n$ meer ontdekt met $n$ groter dan 5. We weten ook niet of ze bestaan. Dat las je misschien al in de bijdrage over Gauss.
	
	Toen Daniël Bernoulli zich in 1733 politiek en academisch beknot voelde onder tsaar Peter de Grote, keerde hij terug naar Bazel en nam Euler zijn leerstoel aan de wiskundefaculteit over. Twee jaar later werd Euler blind aan zijn rechteroog, waarschijnlijk na een infectie door een te lange observatie van de zon. Dat weerhield hem niet om in het jaar 1735 een van zijn roemruchtigste ontdekkingen te doen in zijn publicatie \emph{De summis serierum reciprocarum}. Daarin werkte hij maar liefst drie bewijzen uit van het convergentiegetal uit het probleem van Bazel. In 1747 ontdekte en publiceerde hij er nog een vierde.
	
	Na een verblijf van 13 jaar in Sint-Petersburg kreeg Euler door Frederik de Grote van Pruisen een positie aangeboden in de Academie van Berlijn. Euler ging op dit aanzoek in voor de volgende 24 jaar en werd ondertussen nog halftime doorbetaald door de Academie in Sint-Petersburg. De banden met tsaar Peter de Grote waren nog sterk. In die tijd onderhield Euler een drukke communicatie met Westerse wiskundigen aan andere universiteiten. In deze briefwisseling gaf hij een aanzet tot het bewijs van het, toen nog, vermoeden van Fermat. Dit vermoeden, dat pas in 1994 een stelling geworden is, zegt dat er geen gehele oplossingen $(x,y,z)$ zijn voor de vergelijking $x^n+y^n=z^n$ als $n$ een natuurlijk getal groter dan 2 is. Euler bewees het vermoeden voor $n=3$.
	
	Toen Euler op 59-jarige leeftijd uit de gratie van Frederik de Grote viel, die hem smalend de `wiskundig cyloop' noemde, keerde hij terug naar Rusland. Onder het regime van Catharina de Grote, ging hij onvermoeid verder met zijn publicaties, waardoor zijn naam en faam in Europa ongeëvenaard werden. Kort na zijn terugkeer in Rusland, werd Euler helemaal blind. Zelf vond hij dit geen groot probleem. Een helper - volgens sommige bronnen een kleermaker uit Berlijn zonder wiskundige voorkennis, volgens andere bronnen een zoon van Euler - noteerde wat Euler hem dicteerde. Euler had een fenomenaal geheugen en kende volgens de overlevering de hele Aeneïs van Vergilius uit het hoofd. Hij had dus geen neerslag op papier nodig om omslachtige berekeningen te kunnen uitvoeren. Nadat hij volledig blind geworden was, beweerde Euler dat hij zich veel beter op de wiskunde kon concentreren omdat zijn aandacht niet meer afgeleid kon worden door visuele bijkomstigheden. 
	
	In 1783 overleed Euler aan een hartaanval terwijl hij met zijn kleinzoon speelde. Net ervoor had hij zijn berekeningen afgemaakt over de baan van de pas ontdekte planeet Uranus.    
	
	\subsection*{Convergerende en divergerende reeksen}   
	
	``Als we een oneindig aantal getallen optellen, dan levert dit oneindig op'', meenden de oude Grieken. Dit misverstand leidde tot de gekende paradoxen van Zeno, die van de pijl en die van Achilles en de schildpad. 
	
	Dat de som van oneindig veel positieve getallen toch beperkt kan zijn tot een eindig getal, moet je in de klas zeker aantonen als het over rijen en sommen van rijen gaat. Bij onze leerlingen schiet de intuïtie om dit te kunnen aanvaarden ook soms tekort. 
	
	Het eenvoudigste voorbeeld hiervan gaat over de oneindige rij $1, \frac{1}{2}, \frac{1}{4}, \frac{1}{8}, \dots $ en over de oneindige reeks (of som van een rij) $1+ \frac{1}{2}+ \frac{1}{4}+ \frac{1}{8}+ \dots$ De volgende figuren tonen aan dat deze reeks naar 2 convergeert. Het zijn bewijzen zonder woorden, proofs without words.
	
	\afbeelding[8cm]{img/LoepGeschiedenis_convergentie1.jpg}{Eerste convergentiebewijs}{fig:conv1}
	
	\afbeelding[8cm]{img/LoepGeschiedenis_convergentie2.jpg}{Tweede convergentiebewijs}{fig:conv2}
	
	In het algemeen geven meetkundige rijen met reden of vermenigvuldigingsfactor $q \in ]-1,1[$ aanleiding tot een convergerende reeks. Een nodige voorwaarde voor een rij van positieve getallen om een convergerende reeks te hebben, is het naar nul gaan van de algemene term $t_n$ als $n \to \infty$. Maar deze voorwaarde is uiteraard niet voldoende. De term $t_n$ moet `snel genoeg' naar nul gaan.
	
	Het meest klassieke voorbeeld van een rij met termen $t_n$ die niet snel genoeg naar nul gaan om een convergerende reeks te hebben, is de harmonische rij: 
	
	\[1, \frac{1}{2}, \frac{1}{3}, \frac{1}{4} , \frac{1}{5}, \frac{1}{6}, \frac{1}{7}, \frac{1}{8},\] 
	\[\frac{1}{9}, \frac{1}{10}, \frac{1}{11}, \frac{1}{12}, \frac{1}{13}, \frac{1}{14}, \frac{1}{15}, \frac{1}{16}\dots.\] 
	
	Het klassieke bewijs om aan te tonen dat de bijbehorende reeks divergeert, gaat uit van een \emph{minorante rij} (een rij met kleinere termen dan de gegeven rij) waarvan men makkelijk kan aantonen dat de bijbehorende reeks divergeert. Een geschikte minorante rij is:  
	
	\[1, \frac{1}{2}, \frac{1}{4}, \frac{1}{4} , \frac{1}{8}, \frac{1}{8}, \frac{1}{8}, \frac{1}{8},\]
	\[\frac{1}{16}, \frac{1}{16}, \frac{1}{16}, \frac{1}{16}, \frac{1}{16}, \frac{1}{16}, \frac{1}{16}, \frac{1}{16}\dots.\] 
	
	Als de kleinere (of minorante) reeks naar oneindig gaat dan gaat de grotere (of majorante) reeks zeker naar oneindig. Welnu, door in de minorante rij de termen met gelijke noemers bij elkaar te tellen vinden we oneindig veel keer $\frac{1}{2}$ als som. Dit bewijst dat de minorante reeks naar oneindig gaat. De harmonische reeks gaat dus eveneens naar oneindig. In Kindt (2015) vinden we van dit bewijs ook een meetkundige interpretatie, zie \autoref{fig:div}.
	
	\afbeelding[8cm]{img/LoepGeschiedenis_harmonischereeks.jpg}{Divergentiebewijs van de harmonische reeks}{fig:div}
	
	De harmonische rij bestaat uit de inversen van de termen van de eenvoudigste rekenkundige rij: $1, 2, 3, 4, \dots$. Deze rij heeft een verschil $v=1$ tussen de opeenvolgende termen. 
	
	Ook hier kunnen we een algemener resultaat voorleggen. De rij van de inversen van om het even welke rekenkundige rij, levert een divergerende reeks op, ook al is het verschil $v$ erg groot. Nemen we bijvoorbeeld de rij:
	\[\frac{1}{1001}, \frac{1}{2001},\frac{1}{3001},\frac{1}{4001}, \dots \]
	dan kunnen we de divergentie van de reeks bewijzen door over te gaan naar de minorante rij:
	\[\frac{1}{1001}, \frac{1}{2002},\frac{1}{3003},\frac{1}{4004}, \dots \]
	waarvan we ook weer makkelijk kunnen bewijzen dat de bijbehorende reeks divergeert. Denk hier maar eens over na.
	
	\subsection*{Het probleem van Bazel}
	
	Het convergentieprobleem waarop vooraanstaande wiskundigen  voor 1735 gedurende bijna een eeuw de tanden op hebben stukgebeten is dat van de convergentie van de reeks $\frac{1}{1^2}+\frac{1}{2^2}+\frac{1}{3^2}+\frac{1}{4^2}+ \dots $. 
	
	Dat de reeks convergeert was destijds makkelijk aan te tonen door te werken met majorante rijen (rijen met grotere termen dan de gegeven rij). Indien de bijbehorende reeks van een majorante rij een eindige som heeft, zal de oorspronkelijke rij ook een eindige som hebben. Hiermee kunnen we verder. De gegeven rij:
	\[\frac{1}{1^2}, \frac{1}{2^2},\frac{1}{3^2},\frac{1}{4^2},\frac{1}{5^2},\dots\]
	heeft een majorante rij
	\[\frac{1}{1}, \frac{1}{1\cdot 2},\frac{1}{2\cdot 3},\frac{1}{3\cdot 4},\frac{1}{4\cdot 5},\dots\]
	waarvan we de som kunnen berekenen met een geraffineerde truc. Als we elke breuk splitsen in een verschil van twee breuken, vinden we 2 als som voor de majorante rij:
	
	\begin{align*}
		&\frac{1}{1}+\frac{1}{1\cdot 2}+\frac{1}{2\cdot 3}+\frac{1}{3\cdot 4}+\frac{1}{4\cdot 5}+\dots \\
		=&\frac{1}{1}+\left( \frac{1}{1}-\frac{1}{2} \right)+\left( \frac{1}{2}-\frac{1}{3} \right)+\left( \frac{1}{3}-\frac{1}{4} \right)+\left( \frac{1}{4}-\frac{1}{5} \right)+\dots \\
		=&\frac{1}{1}+ \frac{1}{1}+ \left(-\frac{1}{2} + \frac{1}{2}\right) + \left(-\frac{1}{3} + \frac{1}{3} \right) +\left(-\frac{1}{4} + \frac{1}{4}\right)+\dots \\
		=&\frac{1}{1}+\frac{1}{1} \\
		=&2
	\end{align*}
	
	Bij deze berekening moeten we een zekere nonchalantheid rapporteren. De associativiteit van de optelling bij oneindige reeksen is niet zo evident. Bij convergerende reeksen mogen we hier op steunen. Maar gezien we nog niet aangetoond hebben dat deze reeks convergeert, is de associativiteit hier nattevingerwerk. Euler zelf ging vaak en met succes met de natte vinger te werk.
	
	Uit deze nonchalante berekening besluiten we dat als de reeks $\frac{1}{1^2}+\frac{1}{2^2}+\frac{1}{3^2}+\frac{1}{4^2}+ \dots $ convergeert, ze moet convergeren naar een getal dat kleiner is dan 2. Numeriek kon uitgerekend worden dat de som van de inversen van de kwadratische rij in de buurt van 1,64 ... moest liggen. De wiskundegemeenschap keek bewonderend toe toen Euler op 28-jarige leeftijd bewees dat dit getal precies gelijk was aan $\frac{\pi^2}{6}$. Het opduiken van het getal $\pi$ in een probleem dat niets met cirkels te maken had, was werkelijk onverwacht. Momenteel bestaat er een schat aan bewijzen voor het probleem van Bazel maar ze gaan allemaal verder dan wat we in de klas beogen, zie bijvoorbeeld de drie bewijzen in Aigner (2013) en op het YouTubekanaal 3Blue1Brown (Sanderson 2018).
	
	\subsection*{Andere convergente reeksen}
	
	Een andere rij waarvan Euler op een min of meer intuïtieve manier aantoonde dat de som ervan convergeerde, was de rij van de omgekeerden van de faculteitsgetallen: $\frac{1}{0!}, \frac{1}{1!},\frac{1}{2!},\frac{1}{3!},\dots$. De bijbehorende reeks convergeert naar het getal e, het grondtal van de neperiaanse logaritme. Door middel van deze reeks kon Euler e tot op 23 decimalen nauwkeurig berekenen. Tegelijkertijd met de reeksontwikkeling van het getal e was de reeksontwikkeling voor de exponentiële functie met grondtal e geboren:
	\[e^x=1+x+\frac{x^2}{2!}+\frac{x^3}{3!}+\frac{x^4}{4!}+\cdots.\]
	
	We eindigen hier met een ontdekking uit 1737. Euler bewees dat de rij van de omgekeerden van de priemgetallen een oneindige som had. Je zou dit als volgt kunnen interpreteren: in de hoge regionen zijn de kwadraten schaarser dan de priemgetallen. De rij van de omgekeerden van de kwadraten heeft immers wel een eindige som. Ook dit bewijs van Euler zou nu als onvolledig en te weinig gefundeerd bestempeld worden.
	
	Maken we nog even een overstapje naar de moderne tijd, dan kunnen we opmerken dat de rij van de inversen van de priemtweelingen wel een eindige som heeft. Priemtweelingen zijn priemgetallen die slechts twee eenheden van elkaar verschillen. De kleinste priemtweelingen zijn 2 en 3, 3 en 5, 5 en 7, 11 en 13, 17 en 19 enz \dots. In 1919 liet de Noorse wiskundige Viggo Brun zien dat de  som:
	\[\left(\frac{1}{3}+\frac{1}{5}\right)+\left(\frac{1}{5}+\frac{1}{7}\right)+\left(\frac{1}{11}+\frac{1}{13}\right)+\left(\frac{1}{17}+\frac{1}{19}\right)+\cdots\]
	convergeert naar een constante die we nu de constante van Brun noemen. Er bestaat een raming voor deze constante: $B_2=1,9021...$ maar aangezien de reeks van de inversen van de priemtweelingen tergend traag convergeert, kost het veel moeite om de huidige beste benadering scherper te stellen. Niemand weet momenteel trouwens of het aantal priemtweelingen eindig of oneindig is. Er is dus nog veel onderzoekswerk in de wiskunde te doen in elementaire onderzoeksdomeinen.   
	
\end{multicols}

\stopcontents[inhoudloep]		% om inhoud voor het loep artikel te eindigen

%% BRONNEN %%%%%%%%%%%%%%%%%%%%%%%%%%%%%%%%%%%%%%%%%%%%%%%%%%%%%%%%%%%%%%%%%%%%%
\begin{bronnen}
	\item Aigner, M. \& Ziegler G. (2013). \textit{Raisonements divins. Quelques démonstrations mathématiques particulièrement élégantes}. Springer Verlag France. ISBN 978-2-8178-0399-9, pp 49-55.
	\item Berlinghoff, W. \& Gouvêa F. (2016). \textit{Wortels van de Wiskunde. Een historisch overzicht voor leraren en anderen}. Epsilon Uitgaven. ISBN 978-90-5041-156-1.
	\item Bogomolny, A. (2006). \textit{In the Wasan Spirit}. Eigen website. \url{https://www.cut-the-knot.org/pythagoras/LikelySangaku.shtm}
	\item D’Andrea, C. \& Gómez, E. (2006). \textit{The Broken Spaghetti Noodle.} The American Mathematical Monthly 113(6), pp 555–557. \url{https://doi.org/10.2307/27641981}
	\item De Naeghel, K. (2017). \textit{Wiskunde in actie}. Eigen website. \url{https://koendenaeghel3.webs.com/wiskundeinactie.pdf}
	\item Hautekiet, G., Lefèvre, E., Van Emelen, E. (2014). Beschrijvende statistiek in de tweede graad. \emph{Uitwiskeling 30/3}, pp 44-46.
	\item Huvent, G. (2008). \textit{Sangaku. Le mystère des énigmes géométriques japonaises}. Paris: Dunod. (zie ook Uitwiskeling 26/3.)
	\item \textit{Geometric Probability}. Brilliant.org. \url{https://brilliant.org/wiki/1-dimensional-geometric-probability}
	\item Janssens, D., Roelens, M., Op de Beeck, R. (1989). Historisch geïnspireerde wiskundelessen. \emph{Uitwiskeling 6/1}.
	\item Katz, V.J. (2009). \textit{A history of mathematics, an introduction.} Addison-Wesley.
	\item Kindt, M. (2015). \textit{Wat te bewijzen was}. Freudenthal Instituut.
	\item Launay, M. (2018). \textit{De wetten van de wereld. Het grote verhaal van de wiskunde}. Oorspronkelijke titel: \textit{Le grand roman des maths: de la préhistoire à nos jours}. Uitgeverij Balans. ISBN 978-94-638-2018-9.
	\item Magnello, E. (2010). \textit{Florence Nightingale: the compassionate statistician}. Plus. \url{https://plus.maths.org/content/about-plus}
	\item Meskens, A. (2021). Florence Nightingale, Adolphe Quetelet en het begin van de medische statistiek. \emph{Wiskunde \& Onderwijs 47/187}, pp 5-14.
	\item Moons, F. (2020). Mysterie van de dag: waarom zijn er oneindig veel priemgetallen? \emph{Knack}. \url{https://www.knack.be/nieuws/wetenschap/mysterie-van-de-dag-waarom-zijn-er-oneindig-veel-priemgetallen/}
	\item Nicholaevich, A. (2015). \textit{Moscow will get smarter}. Kolmogorov Livejournal. \url{https://kolmogorov.livejournal.com/90279.html}
	\item \textit{Nightingale's `Coxcombs'}. (z.d.). Understanding Uncertainty. \url{https://understandinguncertainty.org/coxcombs}
	\item \textit{Parabola drawing linkage 2}. (2020). YouTube. \url{www.youtube.com/shorts/zV8Ed1GgKJs}
	\item \textit{Paradoxes in Probability}. (z.d.) Brilliant.org.	\url{https://brilliant.org/wiki/paradoxes-in-probability/#simpsons-paradox}
	\item Richeson, D. (2017). \textit{Tales of Impossibility. The 2000-Year Quest to Solve Mathematical Problems of Antiquity.} Princeton University Press.
	\item Roelens, M., Van den Broeck, L. (2022). Parameterkrommen. \emph{Uitwiskeling 38/2}.
	\item Saidak, F. (2006). A new proof of Euclid's theorem. \emph{The American Mathematical Monthly 113/10}, pp 37-938. \url{https://doi.org/10.2307/27642094}
	
	\item 3blue1brown. (2018). \textit{Why is pi here? And why is it squred? A geometric answer to the Basel Problem}. YouTube. \url{https://www.youtube.com/watch?v=d-o3eB9sfls}
	\item Numberphile \& 3blue1brown. (2021). \textit{Bertrand's paradox (with 3blue1brown)}. YouTube.	\url{https://www.youtube.com/watch?v=mZBwsm6B280}
	\item Struik, D.J. (1994). \textit{Geschiedenis van de wiskunde}. Het Spectrum. ISBN 90-274-2210-9.
	\item Vanden Berghe, G. Devreese, J. (2011). \textit{Wonder en is gheen wonder, de geniale wereld van Simon Stevin}. Uitgeverij Davidsfonds. ISBN 978-90-5826-174-8
	\item Vanden Berghe, G. Viaene, D. Vandamme, L. (2020). \textit{Simon Stevin (1548-1620). Hij veranderde de wereld}. Sterck \& De Vreese. ISBN 978-90-5615-655-8.
	\item van den Bogaart, D. (2020). Het bamboeprobleem. \emph{Euclides. Wortels van de wiskunde} (themanummer).
	\item van Lint, H., Breeman, J. (2015). \textit{Sangaku's. Schoonheid van de meetkunde zonder woorden.} Epsilon Uitgaven. Zebrareeks (42).
	\item Wussing, H. (2010). \textit{Geschiedenis van de wiskunde. Vanaf de wetenschappelijke revolutie tot de twintigste eeuw}. Veen Magazines. ISBN 978-90-8571-218-3.
	\item Young, R. V. (editor). (2015). \textit{Notable Mathematicians, from ancient times to the present}. Gale Research. ISBN 0-7876-3071-3.
\end{bronnen}